% !TeX encoding = UTF-8

% 载入 SJTUThesis 模版
\documentclass[type=bachelor,oneside]{sjtuthesis}
% 选项
%   type=[doctor|master|bachelor],     % 可选（默认：master），论文类型
%   zihao=[-4|5],                      % 可选（默认：-4），正文字号大小
%   lang=[zh|en|de|ja],                % 可选（默认：zh），论文的主要语言
%   review,                            % 可选（默认：关闭），盲审模式
%   [twoside|oneside],                 % 可选（默认：twoside），双页或单页边距模式
%   [openright|openany],               % 可选（默认：openright），奇数页或任意页开始新章
%   math-style=[ISO|TeX],              % 可选 (默认：ISO)，数学符号样式
\usepackage[hidelinks]{hyperref}


% 论文基本配置，加载宏包等全局配置
% !TEX root = ./main.tex

\sjtusetup{
  %
  %******************************
  % 注意：
  %   1. 配置里面不要出现空行
  %   2. 不需要的配置信息可以删除
  %******************************
  %
  % 信息录入
  %
  info = {%
    %
    % 标题
    %
    zh / title           = {基于执行反馈的强化学习代码生成研究},
    en / title           = {Research on Reinforcement Learning for Code Generation Based on Execution Feedback},
    %
    % 标题页标题
    %   可使用“\\”命令手动控制换行
    %
    % zh / display-title   = {上海交通大学学位论文\\ \LaTeX{} 模板示例文档},
    % en / display-title   = {A Sample Document \\ for \LaTeX-based SJTU Thesis Template},
    %
    % 关键词
    %
    zh / keywords        = {代码生成，执行反馈，强化学习，监督微调，PPO},
    en / keywords        = {code generation, execution feedback, reinforcement learning, supervised fine-tuning, PPO},
    %
    % 姓名
    %
    zh / author          = {孙海翔},
    en / author          = {\textbf{Sun Haixiang}},
    %
    % 指导教师
    %
    zh / supervisor      = {高振国},
    en / supervisor      = {\textbf{Gao Zhenguo}},
    %
    % 副指导教师
    %
    % assoc-supervisor  = {某某教授},
    % assoc-supervisor* = {Prof. Uom Uom},
    %
    % 学号
    %
    id              = {522071910038},
    %
    % 学位
    %   本科生不需要填写
    %
    zh / degree          = {学士},
    en / degree          = {Bachelor of Science},
    %
    % 专业
    %
    zh / major           = {数学与应用数学},
    en / major           = {Mathematics and Applied Mathematics},
    %
    % 所属院系
    %
    zh / department      = {数学科学学院},
    en / department      = {School of Mathematical Sciences},
    %
    % 答辩日期
    %   使用 ISO 格式 (yyyy-mm-dd)；默认为当前时间
    %
     date                 = {2026-06},
    %
    % 标题页显示日期
    %   覆盖对应标题页的日期显示，原样输出
    %
     zh / display-date    = {2026 年 5 月},
    %
    % 资助基金
    %
    % zh / fund  = {
    %                {国家 973 项目 (No. 2025CB000000)},
    %                {国家自然科学基金 (No. 81120250000)},
    %              },
    % en / fund  = {
    %                {National Basic Research Program of China (Grant No. 2025CB000000)},
    %                {National Natural Science Foundation of China (Grant No. 81120250000)},
    %              },
  },
  %
  % 风格设置
  %
  style = {%
    %
    % 论文标题页 logo 颜色 (red/blue/black)
    %
    % title-logo-color = black,
  },
  %
  % 名称设置
  %
  name = {
    % bib             = {References},
    % ack             = {谢\hspace{\ccwd}辞},
    % achv            = {攻读学位期间完成的论文},
    digest={},
  },
}

% 使用 BibLaTeX 处理参考文献
%   biblatex-gb7714-2015 常用选项
%     gbnamefmt=lowercase     姓名大小写由输入信息确定
%     gbpub=false             禁用出版信息缺失处理
\usepackage[backend=biber,style=gb7714-2015]{biblatex}
% 文献表字体
% \renewcommand{\bibfont}{\zihao{5}\fixedlineskip{15.6bp}}
% 文献表条目间的间距
\setlength{\bibitemsep}{0pt}
% 导入参考文献数据库
\addbibresource{refs.bib}

% 脚注格式
\usepackage[perpage,bottom,hang]{footmisc}

% 定义图片文件目录与扩展名
\graphicspath{{figures/}}
\DeclareGraphicsExtensions{.pdf,.eps,.png,.jpg,.jpeg}

% 确定浮动对象的位置，可以使用 [H]，强制将浮动对象放到这里（可能效果很差）
% \usepackage{float}

% 固定宽度的表格
% \usepackage{tabularx}

% 使用三线表：toprule，midrule，bottomrule。
\usepackage{booktabs}

% 表格中支持跨行
\usepackage{multirow}

% 表格中数字按小数点对齐
\usepackage{dcolumn}
\newcolumntype{d}[1]{D{.}{.}{#1}}

% 使用长表格
\usepackage{longtable}

% 附带脚注的表格
\usepackage{threeparttable}

% 附带脚注的长表格
\usepackage{threeparttablex}

% 算法环境宏包
\usepackage[ruled,vlined,linesnumbered]{algorithm2e}
% \usepackage{algorithm, algorithmicx, algpseudocode}

% 代码环境宏包
\usepackage{listings}
\lstdefinestyle{lstStyleCode}{%
  aboveskip         = \medskipamount,
  belowskip         = \medskipamount,
  basicstyle        = \ttfamily\zihao{6},
  commentstyle      = \slshape\color{black!60},
  stringstyle       = \color{green!40!black!100},
  keywordstyle      = \bfseries\color{blue!50!black},
  extendedchars     = false,
  upquote           = true,
  tabsize           = 2,
  showstringspaces  = false,
  xleftmargin       = 1em,
  xrightmargin      = 1em,
  breaklines        = false,
  framexleftmargin  = 1em,
  framexrightmargin = 1em,
  backgroundcolor   = \color{gray!10},
  columns           = flexible,
  keepspaces        = true,
  texcl             = true,
  mathescape        = true
}
\lstnewenvironment{codeblock}[1][]{%
  \lstset{style=lstStyleCode,#1}}{}

% 直立体数学符号
\providecommand{\dd}{\mathop{}\!\mathrm{d}}
\providecommand{\ee}{\mathrm{e}}
\providecommand{\ii}{\mathrm{i}}
\providecommand{\jj}{\mathrm{j}}

% 国际单位制宏包
\usepackage{siunitx}

% 定理环境宏包
\usepackage{ntheorem}
% \usepackage{amsthm}

% 绘图宏包
\usepackage{tikz}
\usetikzlibrary{arrows.meta, shapes.geometric}

% 一些文档中用到的 logo
\usepackage{hologo}
\providecommand{\XeTeX}{\hologo{XeTeX}}
\providecommand{\BibLaTeX}{\textsc{Bib}\LaTeX}

% 借用 ltxdoc 里面的几个命令方便写文档
\DeclareRobustCommand\cs[1]{\texttt{\char`\\#1}}
\providecommand\pkg[1]{{\sffamily#1}}

% hyperref 宏包在最后调用
\usepackage{hyperref}

% E-mail
\providecommand{\email}[1]{\href{mailto:#1}{\urlstyle{tt}\nolinkurl{#1}}}


\begin{document}
\setlength{\baselineskip}{20pt}

%TC:ignore

% 标题页
\maketitle

% 原创性声明及使用授权书
\copyrightpage
% 插入外置原创性声明及使用授权书
% 此时必须在导言区使用 \usepackage{pdfpages}
% \copyrightpage[scans/sample-copyright.pdf]

% 前置部分
\frontmatter
{
\fancyhead[LE,RO]{}
{
\ctexset{chapter={afterskip=26bp}}
% 摘要
% !TEX root = ../main.tex

\begin{abstract}[zh]
\addcontentsline{toc}{chapter}{摘 \quad 要}
随着大语言模型和代码大模型的发展，基于自然语言描述自动生成程序代码已经成为软件工程和人工智能交叉领域的重要研究方向。现有代码生成模型通常先经过大规模预训练，再通过监督微调适配特定任务格式。然而，监督微调主要优化参考代码的 token 级似然，与代码生成任务最终关注的功能正确性并不完全一致。对于代码任务而言，生成结果不仅需要语法正确，还需要能够执行并通过测试用例。因此，如何利用代码执行结果构建反馈信号，并进一步优化模型生成策略，是提升代码生成可靠性的关键问题。

本文围绕基于执行反馈的强化学习代码生成方法展开研究，构建了从数据整理、监督微调到 PPO 强化学习优化的完整实验流程。首先，本文基于 MBPP 和 HumanEval 数据集构造统一格式的数据文件，将原始样本整理为 SFT 数据和 RL 数据，并对标准答案进行 runtime 复核，确保测试代码和入口函数可用。其次，本文以 Qwen2.5-Coder-1.5B 作为基座模型，采用 LoRA 进行监督微调，得到 SFT baseline。随后，本文实现了代码生成、代码清洗、入口函数对齐、本地执行测试、奖励计算、优势估计和 PPO 更新等关键模块，形成基于执行反馈的代码生成强化学习闭环。

在奖励函数设计方面，本文比较了原始测试通过率奖励、reward\_v2、reward\_v3 和 reward\_v3b 等多种形式。其中，reward\_v2 将代码可执行性、测试通过比例和全部测试通过奖励结合起来，为 PPO 提供更密集的执行反馈；reward\_v3 和 reward\_v3b 进一步尝试引入 AST 结构相似度和函数签名匹配奖励，用于分析结构奖励的实际效果。实验结果表明，在 MBPP 验证集上，SFT baseline 的 Pass@1 为 0.5714，采用 reward\_v2 的 PPO best 达到 0.6667，追平 Base 模型表现，说明执行反馈 PPO 能够修复小样本 SFT 后模型在验证集上的性能下降。在 HumanEval 测试集上，Base、SFT baseline 和 PPO best 的 Pass@1 分别为 0.5000、0.5244 和 0.5122，表明当前 PPO 的增益主要集中在验证集分布上，跨数据集泛化能力仍有不足。消融实验进一步表明，reward\_v2 是当前最有效的奖励设计，而加入结构相似度的 reward\_v3 和 reward\_v3b 未能提升 Pass@1。

综上，本文验证了基于执行反馈的 PPO 优化在小规模代码生成任务中的可行性，证明执行测试结果可以作为有效奖励信号改善 SFT 后模型的部分生成行为。同时，实验也表明，强代码基座模型本身已经具备较高能力，小规模 SFT 和 PPO 并不必然全面超过基座模型，且当前方法在测试集泛化、奖励设计和 PPO 训练稳定性方面仍存在改进空间。
\end{abstract}

{
\newfontface{\arial}{Arial}[Scale=0.94]
\ctexset{chapter/format+={\arial}}
\begin{abstract}[en]
\addcontentsline{toc}{chapter}{ABSTRACT}
With the rapid development of large language models and code-oriented foundation models, generating program code from natural language descriptions has become an important research topic at the intersection of software engineering and artificial intelligence. Existing code generation models are usually pretrained on large-scale corpora and then adapted to downstream tasks through supervised fine-tuning. However, supervised fine-tuning mainly optimizes token-level likelihood against reference solutions, which is not fully aligned with the functional correctness required by code generation tasks. For generated code, syntactic validity is not sufficient; the code must also be executable and pass the corresponding test cases. Therefore, how to use execution results as feedback signals and further optimize the generation policy is a key problem for improving the reliability of code generation models.

This thesis studies reinforcement learning for code generation based on execution feedback, and builds a complete experimental pipeline from data processing and supervised fine-tuning to PPO-based policy optimization. First, MBPP and HumanEval are converted into unified data formats for SFT and RL training, and runtime validation is conducted to ensure that the reference solutions, entry points, and test code are executable. Second, Qwen2.5-Coder-1.5B is used as the base model, and LoRA-based supervised fine-tuning is applied to obtain the SFT baseline. Then, this thesis implements the key components of an execution-feedback learning loop, including code generation, output cleaning, entry-point alignment, local test execution, reward computation, advantage estimation, and PPO policy update.

For reward design, this thesis compares several reward functions, including the original test-pass-ratio reward, reward\_v2, reward\_v3, and reward\_v3b. reward\_v2 combines executability, test pass ratio, and full-pass bonus to provide denser execution feedback for PPO. reward\_v3 and reward\_v3b further introduce AST structural similarity and function signature matching to examine whether structural rewards can improve code generation performance. Experimental results show that on the MBPP validation set, the SFT baseline achieves a Pass@1 of 0.5714, while the best PPO model with reward\_v2 achieves 0.6667, matching the Base model. This indicates that PPO with execution feedback can recover the performance drop caused by small-scale SFT on the validation distribution. On the HumanEval test set, the Pass@1 scores of Base, SFT baseline, and PPO best are 0.5000, 0.5244, and 0.5122, respectively, showing that the improvement of PPO mainly appears on the validation distribution and does not fully transfer to the cross-dataset test setting. Ablation experiments further show that reward\_v2 is the most effective reward design in this work, while reward\_v3 and reward\_v3b with structural similarity rewards do not improve Pass@1.

In summary, this thesis verifies the feasibility of PPO optimization with execution feedback for small-scale code generation tasks, and shows that execution results can serve as useful reward signals to improve part of the model behavior after SFT. At the same time, the experiments also indicate that a strong code base model already has considerable code generation capability, and small-scale SFT and PPO do not necessarily outperform the base model in all settings. The current method still has limitations in cross-dataset generalization, reward design, and PPO training stability.

\end{abstract}
}

}

{
\ctexset{chapter={afterskip=26bp}}
\renewcommand{\cftchapfont}{\zihao{4}\bfseries}
\renewcommand{\cftsecfont}{\zihao{-4}}
\renewcommand{\cftsubsecfont}{\zihao{5}}
% 目录
\tableofcontents
}
}
% % 插图索引
% \listoffigures*
% % 表格索引
% \listoftables*
% % 算法索引
% \listofalgorithms*
% % 符号对照表
% \input{contents/nomenclature}

%TC:endignore

% 主体部分
\mainmatter

% 正文内容
{
\ctexset{chapter={afterskip=26bp}}
% !TEX root = ../main.tex

\chapter{绪论}

\section{研究背景与意义}

近年来，大语言模型在自然语言理解、文本生成、逻辑推理和代码生成等任务中取得了显著进展。以 Transformer 为代表的神经网络结构通过自注意力机制增强了模型对长距离依赖关系的建模能力，为后续大规模预训练语言模型的发展奠定了基础\cite{Vaswani2017Attention}。在此基础上，研究者开始将大规模预训练方法应用于程序语言建模，使模型能够从海量代码语料中学习编程语言的语法结构、常用 API 调用模式和问题求解范式。与普通自然语言相比，程序语言具有更强的形式化结构和更明确的可执行语义，因此代码生成任务既受到自然语言生成技术发展的推动，也具有区别于普通文本生成的特殊研究价值。

代码生成任务通常要求模型根据自然语言描述、函数签名或上下文代码自动生成满足需求的程序片段。该任务能够应用于自动化编程、智能代码补全、程序修复、测试用例生成和软件工程教育等场景。随着 Codex、Code Llama、DeepSeek-Coder 等代码大模型的发展，模型在 HumanEval、MBPP 等基准测试上的表现不断提升，表明大规模预训练模型已经具备一定的程序合成能力\cite{Chen2021Codex,Roziere2023CodeLlama,Guo2024DeepSeekCoder}。然而，代码生成模型的输出并不总是可靠。自然语言生成任务中可接受的表达差异，在代码任务中可能直接表现为语法错误、运行时异常、入口函数不匹配或边界条件处理错误。即使生成代码在表面上符合题意，也可能无法通过单元测试，因而代码生成任务更强调功能正确性和可执行性。

目前，代码大模型常通过监督微调进一步适配特定任务格式。监督微调能够使模型学习给定数据集中的 prompt-response 格式，并提升模型对任务指令的响应一致性。指令微调和人类反馈训练的成功表明，预训练模型可以通过后训练阶段更好地对齐用户意图\cite{Ouyang2022InstructGPT}。但对于代码生成任务而言，仅依靠标准答案进行监督学习存在明显局限。一方面，代码问题往往存在多种正确实现，模型模仿某一个参考答案并不等价于学会满足测试语义；另一方面，监督微调目标通常基于 token 级别似然优化，而代码质量的关键评价标准是整个程序是否可执行、是否符合函数接口、是否通过测试。也就是说，训练目标与最终评价目标之间存在不一致。

基于执行反馈的强化学习为缓解上述问题提供了可行思路。代码具有天然的可执行性，模型生成的程序可以被解释器、编译器或测试框架直接验证。与人工偏好反馈相比，执行结果能够提供更客观的任务反馈。例如，代码是否可以运行、是否出现异常、通过了多少测试用例、是否全部通过测试，都可以被转化为奖励信号。CodeRL、PPOCoder 和 COMPCODER 等研究已经表明，将编译反馈、单元测试反馈或结构对齐信息引入训练过程，有助于提升模型生成代码的可执行性和功能正确性\cite{Le2022CodeRL,Shojaee2023PPOCoder,Wang2022CompCoder}。因此，围绕执行反馈设计强化学习训练流程，是提升代码生成模型可靠性的一个重要方向。

本课题以“基于执行反馈的强化学习代码生成研究”为题，关注在小规模代码数据条件下，如何构建从监督微调到强化学习优化的完整实验流程。具体而言，本文以 Qwen2.5-Coder-1.5B 作为基座模型，基于 MBPP 和 HumanEval 数据集构建统一格式的数据文件，先通过监督微调获得任务适配模型，再利用模型生成代码后的执行结果构造奖励，并采用 PPO 算法对策略进行进一步优化。本文的研究意义主要体现在以下几个方面。

第一，从训练目标角度看，本文尝试将代码任务的功能正确性纳入模型优化过程。相比仅优化参考答案似然，执行反馈能够直接反映生成程序是否满足测试约束，更接近代码生成任务的最终目标。第二，从工程实现角度看，本文构建了包括数据处理、SFT 训练、代码生成、执行打分、奖励计算、优势估计和 PPO 更新在内的完整闭环，为后续扩展更复杂的代码强化学习方法提供了实验基础。第三，从实验分析角度看，本文不仅比较 Base、SFT 和 PPO 的最终结果，还进一步分析不同 reward 设计、学习率、KL 系数和 rollout 数量对结果的影响，从而更清楚地揭示执行反馈强化学习在小规模代码数据条件下的有效性与局限性。

\section{国内外研究现状}

代码生成研究的发展大致经历了从规则驱动、统计学习、神经网络生成到大语言模型生成的过程。早期程序合成方法更多依赖形式化约束、搜索算法和手工设计的程序空间，而近年来的研究逐渐转向利用大规模神经网络从自然语言和代码语料中学习生成模式。随着大语言模型规模和数据规模不断扩大，代码生成模型已经能够在大量常见编程任务上生成可运行代码，但其功能正确性、鲁棒性和对复杂需求的理解能力仍然存在不足。

\subsection{代码生成模型研究现状}

预训练代码大模型是当前代码生成研究的核心方向之一。Codex 系列工作系统评估了大语言模型在代码生成任务中的能力，并提出 HumanEval 作为衡量函数级程序合成能力的重要基准\cite{Chen2021Codex}。Austin 等人进一步探讨了大语言模型在程序合成任务中的表现，说明大模型可以通过从自然语言描述中学习程序模式完成一定复杂度的代码生成任务\cite{Austin2021ProgramSynthesis}。这些研究推动了代码生成从传统程序合成方法向大规模预训练模型方法转变。

随着开源代码大模型的发展，Code Llama 和 DeepSeek-Coder 等模型进一步提升了开放代码模型的能力边界。Code Llama 在通用语言模型基础上针对代码数据进行训练，覆盖代码补全、代码生成和自然语言到代码等任务\cite{Roziere2023CodeLlama}。DeepSeek-Coder 则强调高质量代码语料和代码智能能力的系统构建，在多种代码任务上展现出较强表现\cite{Guo2024DeepSeekCoder}。这类模型说明，代码大模型已经具备较强的通用编程能力，为下游任务微调和强化学习优化提供了良好基础。

与此同时，研究者也逐渐意识到，单纯生成代码并不意味着代码真正正确。Liu 等人提出 EvalPlus，指出已有 HumanEval 等基准中的测试用例数量有限，可能高估模型生成代码的真实正确率，并通过扩展测试用例更严格地评估模型功能正确性\cite{Liu2023EvalPlus}。Coder-Reviewer reranking 等工作则尝试通过额外的审查模型或重排序机制提升代码生成结果质量\cite{Zhang2022CoderReviewer}。这些研究共同表明，代码生成模型的评价和优化不能只关注文本相似性或模型似然，而应更多关注执行结果和功能正确性。

\subsection{监督微调在代码生成中的应用}

监督微调是将预训练模型适配到特定任务和数据格式的常用方法。在代码生成任务中，监督微调通常将自然语言问题描述、函数签名或上下文代码作为输入，将参考实现作为输出，使模型学习特定数据集中的解题风格和输出格式。对于小规模实验而言，监督微调具有实现简单、训练稳定、结果可解释等优点，因此常被用作强化学习前的热启动阶段。

然而，监督微调也存在目标错位问题。模型在训练时优化的是参考代码的 token 级似然，而测试时评价的是生成代码是否能通过单元测试。由于同一编程问题通常存在多种等价实现，模型不必生成与参考答案完全相同的代码才能正确；反过来，即使生成代码与参考答案在局部 token 上相似，也可能因为边界条件、返回类型或入口函数错误而无法通过测试。因此，监督微调能够提升模型对任务格式的适应能力，却未必稳定提升功能正确性。

此外，小规模监督微调还可能对强基座模型已有能力产生影响。一方面，领域数据能够强化模型在目标任务上的输出模式；另一方面，数据规模过小或分布较窄时，模型也可能过度拟合特定数据形式，从而削弱原本较强的通用代码能力。近期关于领域监督微调的研究也指出，微调对通用能力的影响取决于数据质量、训练设置和任务分布，并不是简单的单向提升或损伤\cite{Lin2025SFTGeneral}。因此，在本文实验中，SFT 不仅作为独立基线，也作为 PPO 训练的初始策略，用于进一步观察执行反馈能否修复 SFT 后的性能下降。

\subsection{强化学习与执行反馈代码生成研究现状}

强化学习在语言模型后训练中的应用主要源于模型输出质量难以完全通过监督数据刻画的问题。InstructGPT 通过人类反馈训练奖励模型，并利用强化学习优化模型输出，使模型更好地遵循用户指令\cite{Ouyang2022InstructGPT}。PPO 作为一种稳定的策略梯度算法，通过限制新旧策略之间的更新幅度，缓解强化学习训练中的策略崩塌问题\cite{Schulman2017PPO}。在语言模型场景中，PPO 常用于根据奖励信号对生成策略进行进一步优化。

与普通自然语言任务相比，代码任务具有更明确的外部反馈来源。CodeRL 将单元测试结果引入代码生成训练过程，通过深度强化学习提升模型在代码生成任务中的表现\cite{Le2022CodeRL}。PPOCoder 进一步提出将预训练程序语言模型与 PPO 结合，利用代码执行结果和结构对齐信息优化代码生成过程，强调代码的可编译性、语法正确性和功能正确性\cite{Shojaee2023PPOCoder}。COMPCODER 则利用编译器反馈提升神经代码生成结果的可编译性，说明编译错误和执行错误均可以作为有效反馈信号\cite{Wang2022CompCoder}。

除强化学习方法外，也有研究从推理和自我修正角度提升代码生成质量。例如，Chain-of-Thought 提示表明大模型可以通过显式中间推理步骤提升复杂问题求解能力\cite{Wei2022CoT}；Self-Debug 相关工作尝试让模型根据执行反馈自我定位和修复错误\cite{Chen2023SelfDebug}。DPO 等偏好优化方法则提供了不显式训练奖励模型的后训练思路\cite{Rafailov2023DPO}。这些工作说明，后训练阶段的反馈利用方式正在从单一监督学习向多种反馈信号和优化目标发展。

总体来看，国内外相关研究已经证明执行反馈对代码生成任务具有重要价值，但仍存在若干值得进一步研究的问题。首先，如何在小规模数据条件下设计稳定有效的奖励函数仍不明确；其次，强化学习更新可能提升当前验证集表现，却不一定带来跨数据集泛化提升；再次，更复杂的结构奖励是否一定优于简单执行奖励，也需要通过消融实验验证。本文正是在这一背景下，围绕 MBPP 和 HumanEval 数据集构建小规模但完整的 SFT+PPO 实验闭环，重点分析执行反馈奖励在代码生成优化中的实际效果。

\section{本文主要工作}

本文围绕基于执行反馈的强化学习代码生成方法，完成了从数据整理、模型训练到实验分析的完整流程。主要工作如下。

\begin{enumerate}
  \item 构建了统一格式的代码生成实验数据。本文基于 MBPP 和 HumanEval 数据集，将原始数据转换为 master、SFT 和 RL 三类格式，并统一字段为任务编号、问题描述、入口函数、参考答案和测试代码等内容。在此基础上，本文进一步进行了本地 runtime 复核，确认标准答案和测试代码可执行，为后续执行反馈训练提供可靠数据基础。
  \item 实现了基于 Qwen2.5-Coder-1.5B 的监督微调流程。本文使用 MBPP 训练数据构造 prompt-response 样本，采用 LoRA 方式进行参数高效微调，获得 SFT baseline。该阶段为后续 PPO 训练提供初始策略，也用于分析小样本 SFT 对强代码基座模型能力的影响。
  \item 构建了代码生成、执行和打分闭环。本文实现了模型生成代码后的清洗、入口函数对齐、本地执行、单元测试统计和 reward 计算流程，支持在无外部沙箱的服务器环境中完成生成代码的自动评测。该流程将不可微的执行结果转化为强化学习可使用的奖励信号。
  \item 设计并比较了多种执行反馈奖励函数。本文从原始测试通过率奖励出发，进一步设计 reward\_v2，将代码可执行、测试通过比例和全部测试通过奖励组合为更密集的反馈信号。同时，本文参考结构奖励思想设计 reward\_v3 和 reward\_v3b，引入 AST 相似度和函数签名匹配，用于分析复杂结构奖励在小数据条件下的实际效果。
  \item 实现并验证了单步 PPO 优化流程。本文在 SFT 模型基础上收集 rollout，计算 reward 和 advantage，并通过 clipped policy loss、value loss 和 KL penalty 更新 LoRA adapter。实验中进一步比较学习率、KL 系数、rollout 数量和 reward 设计对结果的影响，形成较完整的 PPO 调参消融分析。
  \item 完成了验证集、测试集和案例层面的实验分析。本文比较 Base、SFT 和 PPO 在 MBPP 验证集与 HumanEval 测试集上的表现，并分析 PPO 在部分样例中修复 SFT 错误、在部分样例中产生退化的原因，从而更全面地说明方法的有效性和局限性。
\end{enumerate}

\section{本文组织架构}

全文共分为五章，各章安排如下。

第一章为绪论。该章首先介绍代码生成任务的发展背景和研究意义，然后从代码生成模型、监督微调和执行反馈强化学习三个方面综述国内外研究现状，最后概括本文的主要工作和章节安排。

第二章为相关技术与理论基础。该章将介绍代码生成任务、预训练代码大模型、监督微调、强化学习、PPO 算法、执行反馈优化方法以及常用代码生成数据集和评价指标，为后续方法设计奠定理论基础。

第三章为基于执行反馈的强化学习代码生成方法。该章将详细说明本文的数据处理流程、SFT 训练数据构造、代码生成与执行反馈流程、奖励函数设计、PPO 训练流程和系统实现模块。

第四章为实验设计与结果分析。该章将介绍实验环境、数据集统计、SFT 实验、PPO 主实验、超参数消融实验、奖励函数消融实验、HumanEval 泛化分析和典型案例分析，并总结实验结论。

第五章为全文总结与展望。该章将总结本文研究所得主要结论，分析当前工作在数据规模、训练框架、奖励设计和泛化能力方面的局限，并提出后续研究方向。

\section{本章小结}

本章围绕基于执行反馈的强化学习代码生成研究，首先介绍了代码大模型和代码生成任务的发展背景，指出代码生成任务相比普通文本生成更强调可执行性和功能正确性。随后，本章从代码生成模型、监督微调和执行反馈强化学习三个方面梳理了相关研究现状，说明现有方法虽然已经取得显著进展，但仍存在训练目标与功能正确性不一致、执行反馈利用不足以及小规模数据下强化学习稳定性不确定等问题。在此基础上，本章概括了本文的主要研究工作和全文组织结构。后续章节将进一步介绍相关理论基础、方法实现和实验结果分析。

% !TEX root = ../main.tex

\chapter{相关技术与理论基础}

本章介绍本文研究所涉及的相关技术与理论基础。首先，对代码生成任务的定义、输入输出形式和技术难点进行说明；其次，介绍预训练代码大模型和监督微调方法，为本文采用 Qwen2.5-Coder-1.5B 进行 SFT 训练提供背景；然后，对强化学习和 PPO 算法的核心思想进行阐述；最后，介绍基于执行反馈的代码生成优化方法，以及本文实验中使用的数据集和评价指标。

\section{代码生成任务概述}

代码生成是指模型根据给定的自然语言描述、函数签名、上下文代码或接口约束，自动生成满足需求的程序代码。与一般自然语言生成任务不同，代码生成的输出不仅需要在语法层面符合编程语言规范，还需要在语义层面满足题目要求，并能够在给定测试环境中正确运行。对于函数级代码生成任务而言，输入通常包含自然语言问题描述和函数入口信息，输出则是一个完整函数实现。模型生成的代码需要与题目提供的测试用例交互，并返回符合预期的结果。

形式化地，可以将代码生成任务表示为条件生成问题。设输入问题描述为 $x$，目标代码序列为 $y=(y_1,y_2,\ldots,y_T)$，自回归语言模型通过最大化条件概率生成代码：

\begin{equation}
  p_{\theta}(y|x)=\prod_{t=1}^{T}p_{\theta}(y_t|x,y_{<t}),
\end{equation}

其中，$\theta$ 表示模型参数，$y_{<t}$ 表示当前 token 之前已经生成的代码序列。模型在推理阶段根据该条件概率分布逐步生成代码，最终得到完整程序。对于预训练语言模型而言，该建模方式与自然语言生成类似，但代码生成任务具有更强的结构约束。

代码生成任务的难点主要体现在三个方面。第一，代码具有严格的语法约束。模型输出中任何缺失缩进、括号不匹配、变量未定义或函数名错误，都可能导致程序无法执行。第二，代码具有明确的语义约束。某段代码即使语法正确，也可能因为算法逻辑错误、边界条件遗漏或返回类型不符而无法通过测试。第三，代码具有多解性。同一问题往往存在多种正确实现方式，基于文本相似度的评价并不能充分反映程序功能正确性。因此，代码生成评价通常更关注执行结果，而不是仅比较生成代码与参考答案的字面相似度。

本文关注的是函数级 Python 代码生成任务。给定问题描述和入口函数名称，模型需要生成对应函数实现。生成代码会被放入测试环境中执行，并根据测试用例通过情况计算奖励或评价指标。这一任务形式与 HumanEval 和 MBPP 等常用代码生成基准保持一致，也便于将执行结果作为强化学习反馈。

\section{预训练代码大模型}

预训练代码大模型是将大规模语言模型技术应用于程序语言建模的结果。Transformer 架构通过自注意力机制对序列中不同位置的 token 建立关联，使模型能够捕捉长距离依赖关系\cite{Vaswani2017Attention}。在自然语言和代码语料上进行大规模预训练后，模型能够学习词法结构、语法模式、库函数调用、常见算法模板以及自然语言描述与代码之间的对应关系。

从训练目标看，代码大模型通常采用自回归语言建模方式，即给定前文 token 预测下一个 token。对于代码数据而言，这一目标能够促使模型学习代码片段的连续结构。例如，在 Python 函数中，模型需要根据函数签名、变量名和上下文推断后续控制流、表达式和返回值。随着模型规模和训练数据规模扩大，预训练代码模型逐渐具备从自然语言问题描述直接生成代码的能力。

Codex 的研究系统展示了大语言模型在代码生成任务中的能力，并提出 HumanEval 作为重要评价基准\cite{Chen2021Codex}。后续 Code Llama、DeepSeek-Coder 等开源代码模型进一步推动了代码大模型的发展\cite{Roziere2023CodeLlama,Guo2024DeepSeekCoder}。这些模型通常在大规模代码仓库、技术文档和自然语言语料上训练，能够处理代码补全、代码解释、代码翻译和自然语言到代码生成等任务。

尽管预训练代码大模型已经具备较强的通用能力，但在具体任务场景中仍然需要进一步适配。一方面，不同数据集的 prompt 格式、函数签名和输出要求存在差异；另一方面，预训练模型生成代码时可能输出解释性文本、Markdown 代码块或额外测试代码，导致评测脚本难以直接执行。因此，在本文实验中，预训练代码模型既作为 Base 进行直接评测，也作为 SFT 和 PPO 的初始化模型。通过比较 Base、SFT 和 PPO，可以观察监督微调与执行反馈强化学习分别对模型行为产生的影响。

本文选用 Qwen2.5-Coder-1.5B 作为基座模型。该模型属于代码专用因果语言模型，本身已经具备较强的代码生成能力。选择 1.5B 规模模型的原因在于，它能够在单张消费级 GPU 上完成 LoRA 微调和 PPO 实验，兼顾实验可行性和代码生成能力。由于本文研究重点是执行反馈训练流程与奖励设计，而不是从零训练大模型，因此采用强代码基座模型作为实验起点更符合实际应用场景。

\section{监督微调方法}

监督微调是指在预训练模型基础上，使用标注数据进一步训练模型，使其适配特定任务格式和输出分布。在代码生成任务中，监督微调数据通常由输入 prompt 和目标代码 response 构成。模型训练目标是最大化参考答案在输入条件下的似然，即最小化 token 级交叉熵损失：

\begin{equation}
  \mathcal{L}_{\mathrm{SFT}}(\theta)
  = -\sum_{t=1}^{T}\log p_{\theta}(y_t|x,y_{<t}).
\end{equation}

其中，$x$ 表示输入问题描述，$y_t$ 表示参考代码序列中的第 $t$ 个 token。通过该目标，模型学习在给定问题描述后生成参考实现。监督微调的优点是训练过程稳定、实现简单，并且能够快速使模型适应指定输出格式。因此，在许多语言模型后训练流程中，SFT 通常作为强化学习之前的热启动阶段。

在代码生成场景中，SFT 能够带来两个直接作用。第一，它可以使模型更稳定地输出函数实现，而不是混入过多解释性文本。第二，它可以使模型学习目标数据集中的函数命名、解题风格和常见任务形式。对于本文使用的 MBPP 数据而言，SFT 训练样本由自然语言问题描述和标准答案构成，模型通过训练学习从问题描述到函数代码的映射。

但 SFT 也存在局限。由于训练目标是模仿参考答案，模型优化的是 token 级相似性，而不是程序执行正确性。代码问题通常存在多种正确实现方式，某种实现与参考答案 token 差异较大，并不意味着它是错误的；相反，与参考答案局部相似的实现也可能因为边界条件错误而失败。因此，SFT 的优化目标与代码生成最终评价目标之间存在不一致。此外，当 SFT 数据规模较小或分布较窄时，模型可能过度适配训练数据形式，导致部分通用能力下降。本文实验中也将 SFT 作为独立基线，用于分析小样本 SFT 对强代码基座模型的影响。

为了降低训练成本，本文采用 LoRA 进行参数高效微调。LoRA 的核心思想是在冻结原模型大部分参数的情况下，为部分线性变换层引入低秩可训练矩阵，使模型能够通过少量新增参数适配下游任务。设原始权重矩阵为 $W$，LoRA 将参数更新表示为低秩分解：

\begin{equation}
  W' = W + \Delta W = W + BA,
\end{equation}

其中，$A$ 和 $B$ 为可训练低秩矩阵。该方法能够显著降低显存需求和训练参数量，适合在单张 GPU 上完成代码模型微调。本文在 Qwen2.5-Coder-1.5B 的多个注意力和前馈模块中应用 LoRA，以获得可用于后续 PPO 优化的 SFT adapter。

\section{强化学习基础}

强化学习研究智能体如何通过与环境交互学习最优策略。一个典型强化学习问题可以用马尔可夫决策过程表示，包括状态空间 $\mathcal{S}$、动作空间 $\mathcal{A}$、状态转移概率、奖励函数 $r$ 和折扣因子 $\gamma$。在每一步，智能体根据当前状态 $s_t$ 选择动作 $a_t$，环境返回奖励 $r_t$ 并转移到下一个状态 $s_{t+1}$。强化学习的目标是学习策略 $\pi_{\theta}(a|s)$，最大化期望累计回报：

\begin{equation}
  J(\theta)=\mathbb{E}_{\pi_{\theta}}\left[\sum_{t=0}^{T}\gamma^t r_t\right].
\end{equation}

在语言模型生成任务中，可以将输入 prompt 视为状态的一部分，将模型逐步生成 token 的过程视为动作序列。对于代码生成任务，模型完成一次代码生成后，环境通过执行测试返回奖励。因此，代码生成可以被建模为一种序列决策问题：模型策略决定生成什么代码，执行环境根据代码质量给出奖励。

价值函数和优势函数是强化学习中的重要概念。状态价值函数 $V^{\pi}(s)$ 表示在状态 $s$ 下按照策略 $\pi$ 行动所能获得的期望回报；动作价值函数 $Q^{\pi}(s,a)$ 表示在状态 $s$ 下采取动作 $a$ 后继续按照策略 $\pi$ 行动的期望回报。优势函数用于衡量某一动作相对于当前状态平均水平的好坏：

\begin{equation}
  A^{\pi}(s,a)=Q^{\pi}(s,a)-V^{\pi}(s).
\end{equation}

在代码生成实验中，完整程序生成完成后才能执行测试，因此奖励通常是 episode 级别的延迟反馈。本文采用单步 episode 近似，将一次完整代码生成视为一个决策样本，基于执行结果计算 reward 和 advantage。虽然这种实现比完整多步轨迹 PPO 更简化，但足以验证执行反馈能否对代码生成策略产生有效优化。

\section{PPO 算法原理}

PPO 是一种常用的策略梯度强化学习算法，其核心目标是在提升策略收益的同时限制策略更新幅度，避免新策略偏离旧策略过远而导致训练不稳定\cite{Schulman2017PPO}。在语言模型后训练中，PPO 被广泛用于根据奖励模型或外部反馈优化生成策略。

PPO 首先定义新旧策略概率比：

\begin{equation}
  r_t(\theta)=\frac{\pi_{\theta}(a_t|s_t)}{\pi_{\theta_{\mathrm{old}}}(a_t|s_t)}.
\end{equation}

如果直接最大化 $r_t(\theta)A_t$，当优势估计较大时，策略可能发生过度更新。为缓解这一问题，PPO 引入 clipped objective：

\begin{equation}
  \mathcal{L}^{\mathrm{CLIP}}(\theta)
  = \mathbb{E}_t
  \left[
  \min
  \left(
  r_t(\theta)A_t,
  \mathrm{clip}(r_t(\theta),1-\epsilon,1+\epsilon)A_t
  \right)
  \right],
\end{equation}

其中，$\epsilon$ 为裁剪范围超参数。该目标限制新策略相对旧策略的变化幅度，使训练过程更加稳定。

在实际训练中，PPO 通常还会引入值函数损失和 KL 约束。值函数用于估计当前状态的期望回报，从而辅助计算优势；KL 约束则用于限制当前策略与参考策略之间的差异，防止模型为了追求短期奖励而破坏原有语言建模能力。一个常见的总损失形式可以写为：

\begin{equation}
  \mathcal{L}
  = \mathcal{L}_{\mathrm{policy}}
  + c_v\mathcal{L}_{\mathrm{value}}
  + c_{kl}\mathcal{L}_{\mathrm{KL}},
\end{equation}

其中，$c_v$ 和 $c_{kl}$ 分别表示 value loss 和 KL loss 的权重。

本文实现的 PPO 流程包含 clipped policy loss、scalar value head、单步优势估计和 KL penalty。由于本文实验规模较小，PPO 被设计为轻量级单步优化框架：模型先基于 SFT adapter 生成代码，执行环境计算 reward，再用这些 rollout 样本进行一次或多次 mini-batch 更新。该设计能够在有限算力条件下验证执行反馈强化学习的有效性。

\section{基于执行反馈的代码生成优化}

基于执行反馈的代码生成优化利用程序语言可执行的特点，将模型生成结果放入真实或受控执行环境中运行，并根据执行结果构造反馈信号。与自然语言任务中依赖人工偏好或文本相似度不同，代码任务可以直接通过解释器、编译器和单元测试判断输出质量。这使得代码生成任务天然适合引入外部环境反馈。

执行反馈通常包括以下几类信息。第一，可执行性反馈，即代码是否能够被解释器或编译器成功解析和运行。第二，接口一致性反馈，即生成代码是否包含预期入口函数，函数参数和返回形式是否符合要求。第三，测试通过反馈，即生成代码通过了多少单元测试。第四，错误类型反馈，例如语法错误、运行时异常、断言失败或超时。这些反馈既可以用于最终评测，也可以进一步转化为强化学习奖励。

在代码生成强化学习中，奖励函数设计非常关键。如果只在全部测试通过时给予奖励，模型会面对较稀疏的反馈，大量接近正确但未完全通过的代码得不到有效区分。如果奖励过于复杂，例如过度依赖代码结构相似度，又可能使模型优化代理指标而非真正的功能正确性。因此，执行反馈奖励需要在密度和目标一致性之间取得平衡。

本文围绕这一问题设计了多种 reward。原始 reward 主要基于测试通过比例；reward\_v2 在此基础上加入代码可执行奖励和全部通过奖励，使反馈更加密集；reward\_v3 和 reward\_v3b 进一步尝试引入 AST 相似度和函数签名匹配奖励，用于分析结构奖励是否能够提升代码生成质量。第四章实验将表明，在本文小规模数据和单步 PPO 设置下，执行正确性相关奖励比结构相似度奖励更稳定有效。

\section{代码生成数据集与评价指标}

本文实验主要使用 MBPP 和 HumanEval 两个代码生成数据集。MBPP 是面向基础 Python 编程问题的数据集，题目通常由自然语言描述、参考答案和若干测试用例构成。该数据集问题规模相对较小，适合用于监督微调和强化学习 rollout。HumanEval 是函数级代码生成评价基准，题目通常包含函数签名、文档字符串和隐藏测试逻辑，重点考察模型根据问题描述生成正确函数实现的能力\cite{Chen2021Codex}。

在本文实验中，MBPP 被划分为训练集和验证集，用于 SFT 训练、PPO rollout 和验证集评估；HumanEval 作为最终测试集，用于观察模型在跨数据集场景下的泛化能力。为了保证数据可用性，本文将原始数据统一转换为包含 task\_id、prompt、canonical\_solution、entry\_point 和 test\_code 等字段的格式，并对标准答案进行 runtime 复核，确保测试代码能够正确执行。

代码生成任务最常用的指标之一是 Pass@1。对于每个问题，模型生成一个候选答案，如果该答案通过全部测试，则记为通过；Pass@1 表示所有问题中一次生成即通过的比例：

\begin{equation}
  \mathrm{Pass@1}=\frac{N_{\mathrm{pass}}}{N_{\mathrm{total}}},
\end{equation}

其中，$N_{\mathrm{pass}}$ 表示通过全部测试的样本数，$N_{\mathrm{total}}$ 表示总样本数。Pass@1 能够直接反映模型在单次生成场景下的代码正确率，也是本文比较 Base、SFT 和 PPO 的核心指标。

除 Pass@1 外，本文还记录 average reward 和 generation execution failures。average reward 表示所有样本执行反馈奖励的平均值，可以反映模型整体输出质量；generation execution failures 表示生成代码在执行阶段发生失败的数量，可以衡量模型输出的可执行稳定性。对于 PPO 训练过程，还记录 policy loss、value loss、KL loss 和平均策略概率比等指标，用于分析策略更新是否稳定。

\section{本章小结}

本章介绍了本文研究所需的相关技术与理论基础。首先，代码生成任务被定义为一种以自然语言或函数签名为输入、以可执行程序为输出的条件生成问题，其评价重点是功能正确性而非文本相似性。其次，预训练代码大模型为代码生成提供了强大的初始能力，而监督微调能够使模型适配特定任务格式，但也存在训练目标与执行正确性不一致的问题。随后，本章介绍了强化学习和 PPO 算法的基本原理，说明 PPO 通过裁剪目标和 KL 约束限制策略更新幅度。最后，本章讨论了执行反馈在代码生成优化中的作用，并介绍了本文使用的 MBPP、HumanEval、Pass@1 和 reward 等数据集与评价指标。以上内容为第三章提出基于执行反馈的 PPO 代码生成方法奠定了理论基础。

% !TEX root = ../main.tex

\chapter{基于执行反馈的强化学习代码生成方法}

本章介绍本文提出并实现的基于执行反馈的强化学习代码生成方法。整体流程以预训练代码模型为基础，先通过监督微调获得任务适配模型，再利用模型生成代码后的执行结果构造奖励，最后采用 PPO 对模型策略进行进一步优化。与仅依赖参考答案似然的训练方式相比，本文方法将代码是否可执行、是否通过测试、是否满足入口函数约束等信息纳入优化过程，使训练目标更加接近代码生成任务的功能正确性要求。

\section{方法总体框架}

本文方法整体由数据构建、SFT 热启动、代码生成、执行反馈、奖励计算和 PPO 更新六个部分组成。首先，将 MBPP 和 HumanEval 原始数据整理为统一字段格式，并根据训练目标分别构造 SFT 数据和 RL 数据。其次，使用 SFT 数据对 Qwen2.5-Coder-1.5B 进行 LoRA 微调，使模型适配代码生成任务格式。随后，在 PPO 阶段，模型根据 prompt 生成代码，系统对生成结果进行清洗和入口函数对齐，并在本地执行环境中运行测试代码。测试执行结果被转化为 reward 和 advantage，最终用于 PPO 策略更新。

本文方法可以抽象为如下流程：

\begin{enumerate}
  \item 输入代码生成数据集，统一字段并划分 train、valid 和 test。
  \item 构造 SFT 样本，将问题描述转换为 prompt，将标准答案转换为 response。
  \item 使用 LoRA 对基座代码模型进行监督微调，得到 SFT adapter。
  \item 使用当前策略模型生成候选代码，并记录生成文本、token 数量和旧策略 log probability。
  \item 清洗生成代码，对齐入口函数，并在本地执行测试代码。
  \item 根据执行结果计算 reward，并对一个 rollout batch 计算 advantage。
  \item 使用 PPO clipped objective、value loss 和 KL penalty 更新 LoRA adapter。
  \item 对更新后的模型在验证集和测试集上重新执行生成、测试与打分。
\end{enumerate}

从训练目标上看，SFT 阶段主要学习“如何按照给定格式生成代码”，PPO 阶段则进一步学习“什么样的代码更可能通过执行测试”。因此，本文方法不是用 PPO 替代 SFT，而是将 SFT 作为 PPO 的初始化策略。这样既能利用监督学习的稳定性，又能利用执行反馈对功能正确性进行进一步优化。

\section{数据集构建与格式统一}

本文使用 MBPP 和 HumanEval 两个代码生成数据集。MBPP 主要用于训练和验证，HumanEval 用于最终测试。为了使 SFT、PPO 和评测脚本能够共享同一套数据，本文首先将原始数据统一转换为 master 格式。统一后的每条样本包含任务编号、数据来源、问题描述、参考答案、入口函数、测试代码等字段。其核心字段如表 \ref{tab:master-fields} 所示。

\begin{table}[!htp]
  \centering
  \caption{统一 master 数据字段}
  \label{tab:master-fields}
  \begin{tabular}{ll}
    \toprule
    字段名 & 含义 \\
    \midrule
    \texttt{task\_id} & 样本唯一编号 \\
    \texttt{source} & 数据来源，取值为 MBPP 或 HumanEval \\
    \texttt{prompt} & 自然语言题目描述或函数文档字符串 \\
    \texttt{canonical\_solution} & 参考答案代码 \\
    \texttt{entry\_point} & 需要生成的目标函数名 \\
    \texttt{test\_code} & 用于执行验证的测试代码 \\
    \texttt{metadata} & 原始数据中的附加信息 \\
    \bottomrule
  \end{tabular}
\end{table}

在 master 数据基础上，本文进一步构造 SFT 和 RL 两种格式。SFT 数据面向监督微调，主要保留 \texttt{prompt} 和 \texttt{response}，其中 \texttt{response} 来自参考答案。RL 数据面向执行反馈训练，除了 \texttt{prompt} 外，还保留 \texttt{canonical\_solution}、\texttt{entry\_point} 和 \texttt{test\_code}，用于后续执行测试和 reward 计算。

数据构建过程主要由两个脚本完成。\texttt{prepare\_datasets.py} 负责读取原始 MBPP 和 HumanEval 数据，解析参考答案、入口函数和测试代码，并输出统一 master、SFT 和 RL 文件。\texttt{split\_datasets.py} 负责在 MBPP 数据上按照固定随机种子划分训练集和验证集，同时保留 HumanEval 作为测试集。最终数据划分为 MBPP train 385 条、MBPP valid 42 条、HumanEval test 164 条，总计 591 条样本。为了保证执行反馈可靠，本文还对标准答案进行了 runtime 复核，确认所有样本的标准答案在对应测试代码下均可通过。

\section{监督微调数据构造方法}

SFT 阶段的目标是使模型学习代码生成任务的基本输出格式。本文使用如下 prompt 模板构造训练输入：

\begin{codeblock}[language={}]
### Problem:
{prompt}

### Solution:
\end{codeblock}

其中，\texttt{\{prompt\}} 为数据集中原始问题描述。模型需要在该模板后生成对应函数实现。训练样本的 response 为标准答案代码。为了避免模型学习无关文本，SFT 训练时仅对 response 部分计算 loss，而 prompt 部分的 label 被置为忽略值。这样模型不会被优化为复现输入 prompt，而是学习在给定问题后生成解决方案。

本文 SFT 数据格式如表 \ref{tab:sft-fields} 所示。

\begin{table}[!htp]
  \centering
  \caption{SFT 数据字段}
  \label{tab:sft-fields}
  \begin{tabular}{ll}
    \toprule
    字段名 & 含义 \\
    \midrule
    \texttt{task\_id} & 样本编号 \\
    \texttt{source} & 数据来源 \\
    \texttt{prompt} & 输入问题描述 \\
    \texttt{response} & 参考答案代码 \\
    \texttt{entry\_point} & 目标函数名 \\
    \bottomrule
  \end{tabular}
\end{table}

具体实现中，\texttt{train\_sft.py} 定义了 \texttt{SFTJsonlDataset}，负责读取 JSONL 数据并将 prompt 与 response 拼接为完整训练文本。数据整理器 \texttt{DataCollatorForCausalSFT} 对不同长度样本进行 padding，并使用 $-100$ 屏蔽不参与损失计算的位置。模型训练使用 LoRA，仅更新低秩 adapter 参数，从而降低显存占用并提高训练效率。

SFT 阶段输出的 adapter 作为后续 PPO 的初始策略。这样设计的原因是，直接对未适配任务格式的基座模型进行 PPO 更新会导致 rollout 质量较差，执行反馈信号稀疏；而经过 SFT 后，模型已经能够较稳定地生成函数代码，PPO 可以在此基础上进一步优化功能正确性。

\section{代码生成与执行反馈流程}

PPO 阶段的关键是将模型生成的代码转化为可用于强化学习的执行反馈。本文实现了本地代码执行环境 \texttt{LocalCodeSandboxEnv}，用于完成 prompt 到 reward 的转换。该环境并非完整系统沙箱，而是面向本文数据集的轻量本地执行器，负责在受控命名空间中执行生成代码和测试代码。

代码生成首先由 \texttt{CodePolicy} 完成。给定 prompt 后，模型按照 SFT 阶段使用的模板构造输入，并生成最大长度受限的 response。生成结果中可能包含 Markdown 代码块、解释性文本、多余测试代码或重复题目描述，因此不能直接执行。为此，本文设计了 \texttt{clean\_generated\_solution} 函数，对生成文本进行清洗，主要包括：

\begin{enumerate}
  \item 提取 Markdown 代码块中的 Python 代码；
  \item 去除 \texttt{### Problem}、\texttt{### Solution}、\texttt{### Test} 等额外段落；
  \item 截断明显不属于函数实现的解释性文本；
  \item 保留最可能的函数定义代码片段。
\end{enumerate}

由于模型生成的函数名可能与测试代码要求的入口函数不一致，本文进一步进行入口函数对齐。如果生成代码中存在函数定义但函数名与 \texttt{entry\_point} 不一致，系统会尝试将首个函数定义重命名为目标入口函数。这一处理可以减少因函数名不匹配造成的非功能性失败，使评测更关注生成代码的算法逻辑。

执行反馈流程如算法 \ref{alg:execute-score} 所示。

\begin{algorithm}[!htp]
  \caption{代码生成与执行反馈流程}
  \label{alg:execute-score}
  \KwIn{问题描述 $x$，入口函数 $e$，测试代码 $T$，当前策略模型 $\pi_{\theta}$}
  \KwOut{生成代码 $y$，奖励 $r$，执行信息 info}
  根据 prompt 模板构造模型输入\;
  使用 $\pi_{\theta}$ 生成候选代码文本 $y_{\mathrm{raw}}$\;
  对 $y_{\mathrm{raw}}$ 进行代码清洗，得到 $y_{\mathrm{clean}}$\;
  将 $y_{\mathrm{clean}}$ 中的函数名与入口函数 $e$ 对齐\;
  在本地命名空间中执行生成代码\;
  执行测试代码 $T$，统计通过测试数和总测试数\;
  根据执行结果计算 reward 和错误类型\;
  返回生成代码、reward 和执行信息\;
\end{algorithm}

对于 MBPP 数据，测试代码通常由若干 assert 语句组成，系统逐条执行并统计通过数量。对于 HumanEval 数据，测试代码通常包含 \texttt{check(candidate)} 函数，系统会识别该结构并将生成函数作为 candidate 传入。若执行过程中出现异常，则记录错误类型和 traceback，并将该样本标记为执行失败或部分测试失败。

\section{奖励函数设计}

奖励函数是本文方法的核心。本文围绕执行反馈设计了从简单到复杂的多种 reward，用于分析不同反馈信号对 PPO 训练的影响。

\subsection{原始测试通过率奖励}

最直接的奖励设计是基于测试通过比例。设某个样本共有 $N$ 个测试点，生成代码通过了 $N_{\mathrm{pass}}$ 个测试点，则测试通过比例为：

\begin{equation}
  p=\frac{N_{\mathrm{pass}}}{N}.
\end{equation}

原始 reward 可以直接取 $r=p$。该设计简单直观，能够反映代码通过测试的程度。然而，在实际训练中，仅依赖测试通过比例仍然存在反馈稀疏问题。例如，某段代码已经语法正确且入口函数正确，但由于一个边界条件失败而无法全部通过测试，其 reward 与完全不可执行代码之间的区分可能不够明显。因此，本文进一步设计了更密集的 reward\_v2。

\subsection{reward\_v2 密集执行反馈奖励}

reward\_v2 的设计目标是在保持执行正确性导向的同时，为代码生成过程中的中间质量提供更密集的正向反馈。该 reward 由三部分组成：

\begin{equation}
  r_{\mathrm{v2}}
  = 0.1 \cdot I_{\mathrm{exec}}
  + 0.8 \cdot p
  + 0.1 \cdot I_{\mathrm{full}},
\end{equation}

其中，$I_{\mathrm{exec}}$ 表示代码是否可执行，$p$ 表示测试通过比例，$I_{\mathrm{full}}$ 表示是否通过全部测试。最终 reward 被截断到 $[0,1]$ 区间。

该设计体现了代码生成质量的层级结构。首先，代码需要能够被解释器执行，因此可执行性应获得基础奖励。其次，代码通过的测试越多，说明功能越接近正确，应获得更高奖励。最后，全部测试通过代表该样本在当前测试集下功能正确，因此给予额外奖励。相比原始 reward，reward\_v2 能够区分完全不可执行、部分正确和完全正确的代码，更适合作为 PPO 的优化信号。

\subsection{reward\_v3 与 reward\_v3b 结构奖励尝试}

受执行反馈代码生成相关工作的启发，本文进一步尝试将代码结构信息纳入 reward。reward\_v3 在 reward\_v2 思路基础上加入入口函数存在性、AST 结构相似度和函数签名匹配奖励。其设计形式为：

\begin{equation}
  r_{\mathrm{v3}}
  =0.10 I_{\mathrm{exec}}
  +0.10 I_{\mathrm{entry}}
  +0.15 S_{\mathrm{AST}}
  +0.10 S_{\mathrm{sig}}
  +0.45 p
  +0.10 I_{\mathrm{full}},
\end{equation}

其中，$I_{\mathrm{entry}}$ 表示生成代码是否包含目标入口函数，$S_{\mathrm{AST}}$ 表示生成代码与参考答案的 AST 节点类型相似度，$S_{\mathrm{sig}}$ 表示函数参数数量是否与参考答案匹配。AST 相似度通过统计生成代码和参考答案中的 AST 节点类型分布计算，旨在衡量两段代码在结构上的接近程度。

考虑到结构奖励可能过度影响功能正确性目标，本文又设计 reward\_v3b，降低结构奖励权重，提高测试通过比例权重：

\begin{equation}
  r_{\mathrm{v3b}}
  =0.05 I_{\mathrm{exec}}
  +0.05 I_{\mathrm{entry}}
  +0.05 S_{\mathrm{AST}}
  +0.05 S_{\mathrm{sig}}
  +0.70 p
  +0.10 I_{\mathrm{full}}.
\end{equation}

reward\_v3 和 reward\_v3b 的目的并不是直接替代 reward\_v2，而是用于验证结构相似度奖励在本文小数据代码生成场景下是否有效。后续实验表明，结构奖励虽然能改变 average reward，但未能提升 Pass@1，因此本文最终采用 reward\_v2 作为主要奖励函数。

\section{PPO 训练流程设计}

本文 PPO 训练流程由 rollout 收集和 policy update 两个阶段组成。rollout 收集阶段使用当前策略模型生成代码并执行测试，得到包含 prompt、生成代码、reward、旧策略 log probability 和 advantage 的样本。policy update 阶段读取 rollout 文件，计算 PPO loss 并更新 LoRA adapter。

在 rollout 收集阶段，本文使用 \texttt{ppo\_rollout\_loop.py} 完成如下操作：

\begin{enumerate}
  \item 从 RL 数据文件读取样本；
  \item 使用当前 policy 生成代码；
  \item 调用本地执行环境计算 reward；
  \item 记录旧策略下生成 response 的 log probability；
  \item 对一个 batch 内的 reward 减去平均 reward，得到简化 advantage；
  \item 将 rollout 写入 \texttt{rollouts.jsonl}。
\end{enumerate}

早期实现中，advantage 使用 batch reward 均值作为 baseline：

\begin{equation}
  A_i = r_i - \bar{r}.
\end{equation}

后续 PPO update 中进一步引入 scalar value head，对每个样本估计 value，并采用单步 GAE 近似：

\begin{equation}
  R_i = r_i,\quad A_i = r_i - V(s_i).
\end{equation}

由于本文将一次完整代码生成视为单步 episode，因此 return 等于该样本的 reward。该设计简化了多步 trajectory 的复杂度，适合本文在小规模数据和有限算力条件下验证执行反馈优化效果。

在 policy update 阶段，本文使用 \texttt{ppo\_update.py} 对 LoRA adapter 进行更新。对于每条 rollout 样本，脚本重新计算当前策略下 response 的 log probability，并与旧策略 log probability 构造概率比：

\begin{equation}
  \rho_i=\exp(\log \pi_{\theta}(y_i|x_i)-\log \pi_{\theta_{\mathrm{old}}}(y_i|x_i)).
\end{equation}

策略损失采用 PPO 裁剪目标：

\begin{equation}
  \mathcal{L}_{\mathrm{policy}}
  =-\min(\rho_i A_i,\mathrm{clip}(\rho_i,1-\epsilon,1+\epsilon)A_i).
\end{equation}

值函数损失采用均方误差：

\begin{equation}
  \mathcal{L}_{\mathrm{value}}=(V(s_i)-R_i)^2.
\end{equation}

同时，为限制更新后策略偏离参考策略过远，本文引入近似 KL penalty。最终总损失为：

\begin{equation}
  \mathcal{L}
  =
  \mathcal{L}_{\mathrm{policy}}
  +c_v\mathcal{L}_{\mathrm{value}}
  +c_{kl}\mathcal{L}_{\mathrm{KL}}.
\end{equation}

其中，$c_v$ 为 value loss 权重，$c_{kl}$ 为 KL loss 权重。本文实验中主要设置 $c_v=0.5$，$c_{kl}=0.02$，并在消融实验中比较不同学习率、KL 系数和 rollout 数量对结果的影响。

\section{系统实现与关键模块}

本文实现了一个围绕代码生成 SFT 和 PPO 的轻量实验系统。系统主要模块如表 \ref{tab:system-modules} 所示。

\begin{table}[!htp]
  \centering
  \caption{系统关键模块}
  \label{tab:system-modules}
  \begin{tabular}{ll}
    \toprule
    模块或脚本 & 功能 \\
    \midrule
    \texttt{prepare\_datasets.py} & 原始数据读取、字段标准化、测试代码构造 \\
    \texttt{split\_datasets.py} & 数据划分和 SFT/RL 格式生成 \\
    \texttt{train\_sft.py} & LoRA SFT 训练 \\
    \texttt{generate\_execute\_score.py} & 生成代码、执行测试、计算 reward \\
    \texttt{ppo\_rollout\_loop.py} & PPO rollout 收集和 advantage 预处理 \\
    \texttt{ppo\_update.py} & PPO policy update、value head 训练和 adapter 保存 \\
    \texttt{simple\_openhands/rl/bridge.py} & 代码清洗、执行环境、reward 计算和策略封装 \\
    \texttt{simple\_openhands/rl/ppo.py} & value head 和单步 GAE 辅助函数 \\
    \bottomrule
  \end{tabular}
\end{table}

系统的数据流如下。首先，数据处理脚本生成 \texttt{data/final} 目录下的 SFT 和 RL 文件。然后，SFT 训练脚本读取 \texttt{sft\_train.jsonl} 和 \texttt{sft\_valid.jsonl}，输出 LoRA adapter。接着，rollout 脚本加载基座模型和 adapter，在 RL 数据上生成代码并写出 \texttt{rollouts.jsonl}。PPO 更新脚本读取 rollout 文件，更新 adapter 并保存新的 PPO 模型。最后，评测脚本在 valid 和 HumanEval test 上重新生成代码并执行测试，输出 \texttt{summary.json} 和 \texttt{scored\_rollouts.jsonl}。

为了支持论文实验分析，系统还记录了每个样本的任务编号、生成代码、reward、是否通过、通过测试数量、错误类型和 reward 组成项。这些信息不仅用于计算总体指标，也用于后续案例分析。例如，通过比较 Base、SFT 和 PPO 在同一任务上的生成代码，可以判断 PPO 是修复了 SFT 的错误，还是造成了边界条件退化。

需要说明的是，本文实现的 PPO 框架是面向本科论文实验的简化版本。它已经包含 policy loss、value loss、KL penalty、单步 advantage 和执行反馈 reward，但尚未实现完整多步 trajectory、复杂 critic 训练和大规模分布式 rollout。因此，本文更关注在小规模数据和有限算力条件下验证执行反馈强化学习代码生成的可行性，而不是构建工业级大规模 RL 系统。

\section{本章小结}

本章详细介绍了本文提出的基于执行反馈的强化学习代码生成方法。首先，本章给出了整体框架，说明 SFT 热启动和 PPO 执行反馈优化之间的关系。随后，本章介绍了 MBPP 和 HumanEval 数据的字段统一、SFT/RL 数据构造和 runtime 复核流程。接着，本章说明了代码生成后的清洗、入口函数对齐和本地执行反馈流程，并重点介绍了原始 reward、reward\_v2、reward\_v3 和 reward\_v3b 的设计。最后，本章给出了 PPO rollout 收集、单步 advantage 估计、policy update 和系统关键模块。下一章将在此方法基础上，进一步给出实验设置、主实验结果、消融实验和案例分析。

% !TEX root = ../main.tex

\chapter{实验设计与结果分析}

本章围绕第三章提出的基于执行反馈的 PPO 代码生成方法进行实验验证。实验目标包括：第一，比较 Base、SFT 和 PPO 三类模型在验证集和测试集上的表现；第二，分析 PPO 超参数和 reward 设计对结果的影响；第三，通过案例分析说明 PPO 在部分样例中能够修复 SFT 错误，同时也可能造成边界条件退化。实验结果表明，本文实现的 PPO 闭环能够在 MBPP 验证集上改善 SFT 后模型的表现，但在 HumanEval 测试集上的泛化能力仍有限。

\section{实验环境与模型设置}

本文实验主要在 AutoDL 云服务器上完成。服务器配备单张 NVIDIA GeForce RTX 4090 D 显卡，显存为 24GB，CPU 为 16 核，内存为 60GB。软件环境使用 Python 3.12、PyTorch 2.11.0+cu130，并通过 Hugging Face Transformers、PEFT 和 Accelerate 等库完成模型加载、LoRA 微调和生成评测。实验主要工作目录位于 \texttt{/root/autodl-tmp/work/simple\_openhands}，模型缓存目录位于 \texttt{/root/autodl-tmp/models}。

本文使用 Qwen2.5-Coder-1.5B 作为基座模型。该模型是面向代码任务训练的因果语言模型，具备较强的代码生成能力。实验中，Base 表示不加载任何 adapter 的基座模型；SFT baseline 表示在 MBPP 训练集上进行 LoRA 监督微调后的模型；PPO best 表示在 SFT adapter 基础上使用 reward\_v2 和 PPO 更新得到的最佳模型。

SFT 训练的主要配置如表 \ref{tab:sft-config} 所示。

\begin{table}[!htp]
  \centering
  \caption{SFT 训练主要配置}
  \label{tab:sft-config}
  \begin{tabular}{ll}
    \toprule
    参数 & 取值 \\
    \midrule
    基座模型 & Qwen2.5-Coder-1.5B \\
    训练文件 & \texttt{data/final/sft\_train.jsonl} \\
    验证文件 & \texttt{data/final/sft\_valid.jsonl} \\
    训练轮数 & 5 \\
    最大长度 & 1024 \\
    batch size & 2 \\
    gradient accumulation steps & 8 \\
    learning rate & $1\times10^{-4}$ \\
    LoRA rank & 16 \\
    LoRA alpha & 32 \\
    LoRA dropout & 0.1 \\
    target modules & q\_proj, k\_proj, v\_proj, o\_proj, gate\_proj, up\_proj, down\_proj \\
    \bottomrule
  \end{tabular}
\end{table}

PPO 训练采用单步 episode 形式，主要配置如表 \ref{tab:ppo-config} 所示。由于实验算力和数据规模有限，本文没有进行大规模多轮 PPO，而是重点分析不同超参数和 reward 设计对一次 PPO update 的影响。

\begin{table}[!htp]
  \centering
  \caption{PPO 训练主要配置}
  \label{tab:ppo-config}
  \begin{tabular}{ll}
    \toprule
    参数 & 取值 \\
    \midrule
    初始策略 & SFT adapter \\
    epochs & 1 \\
    batch size & 4 \\
    clip epsilon & 0.2 \\
    value loss coefficient & 0.5 \\
    gamma & 1.0 \\
    GAE lambda & 1.0 \\
    learning rate & $1\times10^{-5}$ 或 $5\times10^{-6}$ \\
    KL coefficient & 0.02 或 0.01 \\
    rollout 数量 & 42 或 100 \\
    主要 reward & reward\_v2 \\
    \bottomrule
  \end{tabular}
\end{table}

\section{数据集统计与可用性验证}

本文实验使用 MBPP 和 HumanEval 两个数据集。MBPP 用于 SFT 训练、SFT 验证、PPO rollout 和 PPO 验证；HumanEval 作为最终测试集，用于评价模型跨数据集泛化能力。最终数据划分如表 \ref{tab:dataset-stat} 所示。

\begin{table}[!htp]
  \centering
  \caption{最终数据集统计}
  \label{tab:dataset-stat}
  \begin{tabular}{llll}
    \toprule
    文件 & 样本数 & 来源 & 用途 \\
    \midrule
    \texttt{sft\_train.jsonl} & 385 & MBPP & SFT 训练 \\
    \texttt{sft\_valid.jsonl} & 42 & MBPP & SFT 验证 \\
    \texttt{rl\_train.jsonl} & 385 & MBPP & PPO rollout 训练池 \\
    \texttt{rl\_valid.jsonl} & 42 & MBPP & PPO 验证集 \\
    \texttt{rl\_test\_humaneval.jsonl} & 164 & HumanEval & 最终测试集 \\
    \bottomrule
  \end{tabular}
\end{table}

为了保证执行反馈的可靠性，本文对数据进行了 runtime 复核。复核过程将标准答案与测试代码放入执行环境中运行，确认测试代码可以正确调用入口函数并完成断言检查。复核结果如表 \ref{tab:runtime-check} 所示。

\begin{table}[!htp]
  \centering
  \caption{数据 runtime 可用性复核结果}
  \label{tab:runtime-check}
  \begin{tabular}{lrrrr}
    \toprule
    复核集 & 输入样本数 & 通过样本数 & 失败样本数 & pass rate \\
    \midrule
    全量格式复核 & 591 & 591 & 0 & 1.0 \\
    MBPP valid runtime 复核 & 42 & 42 & 0 & 1.0 \\
    MBPP train runtime 复核 & 385 & 385 & 0 & 1.0 \\
    HumanEval runtime 复核 & 164 & 164 & 0 & 1.0 \\
    \bottomrule
  \end{tabular}
\end{table}

该结果说明，本文使用的数据字段、入口函数和测试代码均可被当前执行环境正确处理。由于本文的 reward 直接依赖测试执行结果，数据可执行性复核是后续实验可信性的基础。

\section{SFT 实验设置与结果}

SFT 阶段使用 MBPP train 的 385 条样本进行训练，验证集为 MBPP valid 的 42 条样本。模型采用 LoRA 参数高效微调，训练 5 个 epoch。训练完成后，SFT adapter 的验证集 \texttt{eval\_loss} 为 0.7472，训练耗时约 5 分钟。

为了评价 SFT 是否提升代码生成能力，本文将 Base 和 SFT baseline 在验证集与 HumanEval 测试集上进行对比。结果显示，在 MBPP valid 上，Base 模型的 Pass@1 为 0.6667，而 SFT baseline 为 0.5714；在 HumanEval test 上，Base 的 Pass@1 为 0.5000，而 SFT baseline 为 0.5244。该结果说明，小规模 SFT 并未在所有数据分布上稳定提升模型能力。

这一现象有两点解释。首先，Qwen2.5-Coder-1.5B 本身是较强代码基座模型，在统一清洗和后处理口径下已经具备较高 valid 表现。其次，本文 SFT 数据规模较小，且主要来自 MBPP，模型可能学习到特定数据格式和解题风格，但不一定能全面提升泛化能力。因此，本文后续 PPO 实验更关注执行反馈能否修复 SFT 后在验证集上的性能下降。

\section{PPO 主实验结果}

本文首先比较 Base、SFT baseline 和 PPO best 在 MBPP valid 上的结果。PPO best 使用 reward\_v2、42 条 rollout、learning rate 为 $1\times10^{-5}$、KL 系数为 0.02。验证集结果如表 \ref{tab:valid-main} 所示。

\begin{table}[!htp]
  \centering
  \caption{rl\_valid 上 Base、SFT 与 PPO best 对比}
  \label{tab:valid-main}
  \begin{tabular}{lrrrl}
    \toprule
    模型 & Pass@1 & full pass & average reward & 结论 \\
    \midrule
    Base & 0.6667 & 28/42 & 0.7381 & 强基座模型在 valid 上已较强 \\
    SFT baseline & 0.5714 & 24/42 & 0.6524 & 小样本 SFT 后 valid 表现下降 \\
    PPO best & 0.6667 & 28/42 & 0.7190 & 相比 SFT 明显提升，追平 Base \\
    \bottomrule
  \end{tabular}
\end{table}

从表 \ref{tab:valid-main} 可以看出，PPO best 将 Pass@1 从 SFT baseline 的 0.5714 提升到 0.6667，full pass 样本数从 24 个提升到 28 个。这说明执行反馈 PPO 能够在当前验证集分布上修复 SFT 后模型的性能下降。需要注意的是，PPO best 的 average reward 略低于 Base，但 Pass@1 与 Base 相同。这表明 Base 在部分未全通过样本上可能获得更高部分分，而 PPO best 更集中地提升了全通过样本数量。

\section{PPO 超参数消融实验}

为了分析 PPO 训练稳定性，本文围绕 learning rate、KL 系数和 rollout 数量进行了多组消融实验。所有实验均从 SFT adapter 出发，区别在于 PPO update 的超参数和 rollout 来源。

\subsection{learning rate 对结果的影响}

在原始 reward 条件下，exp1 使用 learning rate 为 $1\times10^{-5}$，Pass@1 达到 0.6190；exp3 将 learning rate 降低到 $5\times10^{-6}$，Pass@1 为 0.5952。虽然 exp3 的 mean ratio 更接近 1，说明策略更新更保守，但最终效果不如 exp1。这说明在 42 条 rollout 的小规模设置下，过度降低学习率会削弱 PPO 对执行反馈的利用能力。

在 reward\_v2 条件下也观察到类似现象。exp1\_reward\_v2 使用 $1\times10^{-5}$，Pass@1 达到 0.6667；exp2\_reward\_v2\_lr5e6 使用 $5\times10^{-6}$，Pass@1 为 0.6190。由此可见，在本文实验设置下，$1\times10^{-5}$ 更适合作为 PPO 更新学习率。

\subsection{KL 系数对结果的影响}

KL 系数用于限制更新后策略偏离参考策略。exp1 使用 KL 系数 0.02，Pass@1 为 0.6190；exp2 将 KL 系数降低到 0.01，Pass@1 降至 0.5476，且 mean ratio 接近 2。该结果说明，过弱的 KL 约束会导致策略偏移过大，使模型偏离 SFT 后已有的代码生成能力。

在语言模型 PPO 中，KL 约束的作用不仅是保持训练稳定，也是在一定程度上保留基座模型和 SFT 模型中已有的通用语言建模能力。对于本文这种小规模代码数据场景，如果 KL 约束过弱，模型可能过度拟合少量 rollout 中的奖励信号，导致整体性能下降。因此，本文最终选择 KL 系数 0.02 作为主要设置。

\subsection{rollout 数量对结果的影响}

本文还比较了 42 条 rollout 与 100 条 rollout 的效果。直观上，更多 rollout 可能提供更丰富的反馈信号，但实验结果并未支持这一点。exp4 使用 rl\_train 中抽取的 100 条 rollout，learning rate 为 $1\times10^{-5}$，Pass@1 为 0.5714；exp5 在 100 条 rollout 基础上将 learning rate 降低到 $5\times10^{-6}$，Pass@1 仍为 0.5714。

这一结果表明，rollout 数量并非越大越好。本文 PPO 实现属于单步近似版本，value 估计和 advantage 估计都较简单。当 rollout 数量扩大时，样本分布更复杂，但 update 机制没有同步增强，更多样本未必能转化为更稳定的优化收益。相比之下，42 条 rollout 与当前 update 强度更加匹配。

完整 PPO 调参结果如表 \ref{tab:ppo-ablation} 所示。

\begin{table}[!htp]
  \centering
  \caption{PPO 超参数消融结果}
  \label{tab:ppo-ablation}
  \begin{tabular}{llllrrr}
    \toprule
    实验组 & rollout & reward & lr & KL & Pass@1 & avg reward \\
    \midrule
    exp1 & 42 & 原始 & $1e^{-5}$ & 0.02 & 0.6190 & 0.6270 \\
    exp2 & 42 & 原始 & $1e^{-5}$ & 0.01 & 0.5476 & 0.5794 \\
    exp3 & 42 & 原始 & $5e^{-6}$ & 0.02 & 0.5952 & 0.6190 \\
    exp4 & 100 & 原始 & $1e^{-5}$ & 0.02 & 0.5714 & 0.6111 \\
    exp5 & 100 & 原始 & $5e^{-6}$ & 0.02 & 0.5714 & 0.6032 \\
    \bottomrule
  \end{tabular}
\end{table}

\section{奖励函数消融实验}

奖励函数是影响 PPO 效果的关键因素。本文比较了原始 reward、reward\_v2、reward\_v3 和 reward\_v3b。原始 reward 主要依赖测试通过比例；reward\_v2 增加代码可执行奖励和全通过奖励；reward\_v3 和 reward\_v3b 进一步加入 AST 相似度和函数签名匹配等结构奖励。

奖励函数消融结果如表 \ref{tab:reward-ablation} 所示。

\begin{table}[!htp]
  \centering
  \caption{奖励函数消融结果}
  \label{tab:reward-ablation}
  \begin{tabular}{llrrl}
    \toprule
    实验组 & reward 设置 & Pass@1 & avg reward & 结论 \\
    \midrule
    SFT baseline & reward\_v2 评测口径 & 0.5714 & 0.6524 & 新 reward 口径基线 \\
    exp1 & 原始 reward & 0.6190 & 0.6270 & 首次优于 SFT \\
    exp1\_reward\_v2 & 执行密集奖励 & 0.6667 & 0.7190 & 当前最佳 \\
    exp2\_reward\_v2 & reward\_v2 + 小学习率 & 0.6190 & 0.6762 & 不如 exp1\_reward\_v2 \\
    exp1\_reward\_v3 & 加入结构奖励 & 0.5952 & 0.7311 & reward 高但 Pass@1 低 \\
    exp1\_reward\_v3b & 降低结构奖励权重 & 0.5952 & 0.6749 & 未改善 Pass@1 \\
    \bottomrule
  \end{tabular}
\end{table}

实验表明，reward\_v2 是本文最有效的奖励设计。相比原始 reward，reward\_v2 更密集，能够区分代码可执行、部分测试通过和全部测试通过等不同层级的输出质量。使用 reward\_v2 后，PPO 在 valid 上的 Pass@1 从 SFT 的 0.5714 提升到 0.6667。

reward\_v3 和 reward\_v3b 的结果则说明，奖励函数并非越复杂越好。reward\_v3 的 average reward 达到 0.7311，高于 reward\_v2 的 0.7190，但 Pass@1 仅为 0.5952。这表明 AST 相似度和函数签名匹配等结构奖励会提高代理 reward 数值，却未必提升最终功能正确性。reward\_v3b 降低结构奖励权重后，Pass@1 仍未改善。因此，本文最终选择 reward\_v2 作为主奖励函数。

\section{HumanEval 测试集泛化分析}

为了验证模型是否具备跨数据集泛化能力，本文在 HumanEval 测试集上对 Base、SFT baseline 和 PPO best 进行最终评测。结果如表 \ref{tab:humaneval-main} 所示。

\begin{table}[!htp]
  \centering
  \caption{HumanEval 测试集结果}
  \label{tab:humaneval-main}
  \begin{tabular}{lrrrr}
    \toprule
    模型 & Pass@1 & full pass & avg reward & execution failures \\
    \midrule
    Base & 0.5000 & 82/164 & 0.5457 & 7 \\
    SFT baseline & 0.5244 & 86/164 & 0.5707 & 2 \\
    PPO best & 0.5122 & 84/164 & 0.5610 & 0 \\
    \bottomrule
  \end{tabular}
\end{table}

从表 \ref{tab:humaneval-main} 可以看出，SFT baseline 在 HumanEval 上略优于 Base，Pass@1 从 0.5000 提升到 0.5244；PPO best 的 Pass@1 为 0.5122，略低于 SFT baseline，但高于 Base。同时，PPO best 的 execution failures 为 0，说明经过执行反馈优化后，模型输出在可执行稳定性上有所改善。

该结果说明，PPO 的主要收益集中在 MBPP valid 分布上，并未稳定迁移到 HumanEval 测试集。造成这一现象的原因可能包括：第一，PPO rollout 主要来自 MBPP 分布，模型优化方向更贴近 MBPP 题型；第二，本文 PPO update 规模较小，无法充分学习跨数据集泛化能力；第三，reward\_v2 虽然提高了执行反馈密度，但仍主要依赖当前测试代码，可能对验证集分布产生局部适配。

\section{案例分析}

为了进一步解释实验结果，本文选取典型样例比较 Base、SFT 和 PPO 的输出差异。案例分析可以更直观地说明 PPO 的作用：一方面，PPO 能够修复 SFT 的部分错误；另一方面，PPO 也可能损伤原本正确的边界条件处理。

\subsection{PPO 修复 SFT 错误案例}

第一个正向案例来自 HumanEval/37，任务要求对列表中偶数下标位置的元素排序，而奇数下标位置保持不变。Base 输出混入了额外测试和解释性文本，SFT 输出虽然尝试提取偶数下标元素，但使用了未定义变量 $i$，导致代码无法正确执行。PPO 输出如下：

\begin{codeblock}[language=Python]
def sort_even(l: list):
    l_even = sorted(l[::2])
    return [l_even[i//2] if i%2==0 else l[i] for i in range(len(l))]
\end{codeblock}

该实现显式构造偶数下标元素的排序列表，并通过 $i//2$ 将排序后的偶数位元素放回原位置，同时保留奇数下标元素不变，因此通过测试。该案例说明，执行反馈能够引导 PPO 修复 SFT 中的索引逻辑错误。

第二个正向案例来自 HumanEval/67，任务要求从字符串中提取苹果和橘子数量，并用总水果数减去二者得到芒果数。SFT 输出使用 \texttt{s[0]} 和 \texttt{s[-1]} 拼接数字，只能处理个位数情形；PPO 输出如下：

\begin{codeblock}[language=Python]
def fruit_distribution(s,n):
  return n - sum([int(i) for i in s.split() if i.isdigit()])
\end{codeblock}

该实现通过 \texttt{split()} 提取字符串中的数字并求和，更符合题意，也能处理多位数数量。该案例说明，PPO 在部分任务上能够减少 SFT 的硬编码倾向，提高实现的泛化性。

\subsection{PPO 退化案例}

负向案例来自 HumanEval/131，任务要求返回整数中所有奇数位数字的乘积，如果不存在奇数位数字则返回 0。Base 和 SFT 都显式维护 \texttt{has\_odd} 标志，因此能正确处理全偶数输入。PPO 输出如下：

\begin{codeblock}[language=Python]
def digits(n):
  return reduce(lambda x,y: x*y, (int(i) for i in str(n) if int(i)%2!=0), 1)
\end{codeblock}

该实现能够计算存在奇数位数字时的乘积，但当输入所有数字均为偶数时，生成器为空，\texttt{reduce} 会返回初始值 1，而题目要求返回 0。因此，PPO 在该样例上损伤了 SFT 原有的边界条件处理能力。

该案例说明，PPO 并不总是带来提升。由于本文使用的是小规模 rollout 和单步 PPO 更新，模型可能在部分样例上为了提升局部 reward 而牺牲原有边界条件能力。这也解释了为什么 PPO best 在 valid 上能够追平 Base，但在 HumanEval 测试集上略低于 SFT baseline。

\subsection{方法能力边界分析}

综合案例可以看出，执行反馈对代码生成优化具有明确价值，但也存在边界。对于索引逻辑、字符串数字提取等可以被测试用例直接约束的问题，PPO 有机会通过 reward 信号修复 SFT 输出错误；但对于测试覆盖不足或边界条件复杂的问题，PPO 可能学习到通过当前样本的短路径，而不是获得更稳健的通用算法。

因此，本文方法更适合被理解为一种面向当前任务分布的策略调整方法，而不是一次训练即可全面提升模型通用编程能力的方法。若要进一步提升泛化能力，需要扩大 rollout 数据规模，引入更多分布多样的测试样本，并设计更稳健的 reward 与 critic 训练机制。

\section{实验结论}

综合以上实验，可以得到以下结论。

第一，本文实现的 SFT+PPO 训练闭环是可行的。从数据处理、SFT 训练、rollout 收集、执行测试、reward 计算到 PPO update，各模块均已完成并通过实验验证。

第二，在 MBPP valid 上，PPO best 将 SFT baseline 的 Pass@1 从 0.5714 提升到 0.6667，说明执行反馈 PPO 能够修复小样本 SFT 后模型在当前分布上的性能下降。

第三，PPO 对超参数和 reward 设计较敏感。单独降低 KL 系数会导致策略偏移过大，单独降低学习率会削弱收益，扩大 rollout 数量也没有自动提升效果。当前最佳结果来自 42 条 rollout、learning rate 为 $1\times10^{-5}$、KL 系数为 0.02 和 reward\_v2 的组合。

第四，reward\_v2 是当前最有效的奖励函数。相比原始 reward，它提供了更密集的执行反馈；相比 reward\_v3 和 reward\_v3b，它更直接地对齐测试通过率和功能正确性。

第五，当前方法的泛化能力仍有限。PPO best 在 HumanEval 上略低于 SFT baseline，说明当前 PPO 更偏向验证集分布优化，尚未稳定提升跨数据集测试表现。

\section{本章小结}

本章对本文方法进行了实验验证和结果分析。首先，本章介绍了实验环境、模型设置、数据集统计和 runtime 复核结果。随后，本章比较了 Base、SFT 和 PPO best 的主实验结果，证明 reward\_v2 条件下的 PPO 能够在 MBPP valid 上显著改善 SFT baseline。接着，本章通过超参数消融和奖励函数消融说明 PPO 效果依赖于合理的更新强度和 reward 设计。最后，本章通过 HumanEval 泛化分析和案例分析指出，当前方法仍存在跨数据集泛化不足和边界条件退化问题。下一章将总结全文工作，并讨论后续改进方向。

% !TEX root = ../main.tex

\chapter{全文总结与展望}

\section{主要结论}

本文围绕“基于执行反馈的强化学习代码生成研究”这一主题，构建了从数据整理、监督微调到 PPO 强化学习优化的完整实验流程，并在 MBPP 和 HumanEval 数据集上进行了实验验证。本文的核心思想是利用代码任务天然具备的可执行特性，将模型生成代码后的运行结果、测试通过情况和错误信息转化为奖励信号，从而使模型优化目标更接近代码生成任务的功能正确性要求。

本文首先对 MBPP 和 HumanEval 数据进行了统一处理，将原始数据转换为适用于 SFT 和 PPO 的两类格式。SFT 数据主要包含 prompt 和 response，用于监督微调；RL 数据保留 prompt、参考答案、入口函数和测试代码，用于生成、执行和 reward 计算。最终实验数据包含 MBPP train 385 条、MBPP valid 42 条和 HumanEval test 164 条。为了保证执行反馈可靠，本文对标准答案进行了 runtime 复核，确认所有样本的测试代码和入口函数均可被执行环境正确处理。

在模型训练方面，本文以 Qwen2.5-Coder-1.5B 作为基座模型，采用 LoRA 进行监督微调，得到 SFT baseline。随后，本文构建了代码生成、执行测试、reward 计算、advantage 估计和 PPO policy update 的闭环流程。该流程能够读取 RL 格式数据，使用当前策略模型生成代码，对生成结果进行清洗和入口函数对齐，在本地执行环境中运行测试，并根据执行反馈计算奖励。PPO 更新阶段则通过 clipped policy loss、value loss 和 KL penalty 对 LoRA adapter 进行进一步优化。

实验结果表明，本文实现的 PPO 方法能够在 MBPP valid 上改善 SFT 后模型的表现。在统一评测口径下，SFT baseline 在 rl\_valid 上的 Pass@1 为 0.5714，PPO best 在 reward\_v2 设置下达到 0.6667，追平 Base 模型的 Pass@1。该结果说明，在当前验证集分布下，基于执行反馈的 PPO 能够修复小样本 SFT 后模型出现的性能下降，使模型生成代码更符合测试执行要求。

同时，实验也表明，Qwen2.5-Coder-1.5B 本身已经具备较强的代码生成能力。Base 模型在 rl\_valid 上的 Pass@1 已经达到 0.6667，在 HumanEval 测试集上的 Pass@1 为 0.5000。这说明小规模 SFT 和 PPO 并不必然全面超过强代码基座模型，而更适合作为面向特定任务分布的策略调整方法。本文实验中，SFT 在 HumanEval 上相对 Base 有小幅提升，Pass@1 从 0.5000 提升到 0.5244；PPO best 在 HumanEval 上的 Pass@1 为 0.5122，略低于 SFT baseline，但 execution failures 降为 0。该结果说明 PPO 提升主要集中在验证集分布上，跨数据集泛化能力仍有限。

在 reward 设计方面，本文比较了原始测试通过率奖励、reward\_v2、reward\_v3 和 reward\_v3b。实验表明，reward\_v2 是当前最有效的奖励函数。该奖励将代码可执行、测试通过比例和全部测试通过三种信号组合起来，使 PPO 能够获得比原始 reward 更密集的执行反馈。在 reward\_v2 下，PPO best 的 valid Pass@1 达到 0.6667。相比之下，reward\_v3 和 reward\_v3b 虽然引入了 AST 相似度和函数签名匹配等结构奖励，但 Pass@1 均为 0.5952，未能超过 reward\_v2。这说明在本文小数据和单步 PPO 设置下，结构相似性奖励与最终功能正确性之间存在偏差，执行测试反馈仍是更可靠的优化信号。

在 PPO 调参方面，本文比较了 learning rate、KL 系数和 rollout 数量的影响。实验发现，单独降低 KL 系数会导致策略偏移过大，使模型表现退化；单独降低 learning rate 可以使更新更保守，但收益也随之减弱；将 rollout 数量从 42 扩大到 100 并没有自动带来性能提升。当前最佳设置为 42 条 rollout、learning rate 为 $1\times10^{-5}$、KL 系数为 0.02、batch size 为 4，并使用 reward\_v2。该结果说明，PPO 效果依赖于 reward 密度、rollout 规模和 update 强度之间的匹配。

综合来看，本文的主要结论可以概括为以下几点：第一，执行反馈可以有效补充 SFT 的 token 级训练目标，使代码生成优化更接近功能正确性；第二，在小规模数据条件下，PPO 能够在验证集分布上改善 SFT 后模型表现，但不一定稳定提升跨数据集泛化能力；第三，奖励函数设计是影响 PPO 效果的关键因素，简单、直接且与测试通过率强相关的 reward\_v2 比复杂结构奖励更有效；第四，PPO 训练需要合适的 KL 约束和更新强度，更多 rollout 或更复杂 reward 并不必然带来更好结果。

\section{本文工作的局限性}

尽管本文完成了从 SFT 到 PPO 的完整代码生成强化学习实验闭环，并取得了一定实验结果，但当前工作仍存在若干局限。

首先，本文使用的数据规模较小。SFT 和 PPO 主要基于 MBPP 数据，其中训练样本为 385 条，验证样本为 42 条。虽然这些数据已经足以验证方法流程和 reward 设计，但相对于代码大模型训练和强化学习后训练所需的数据规模仍然较小。小规模数据容易导致模型对特定题型和测试风格产生局部适配，限制其跨数据集泛化能力。HumanEval 测试结果也表明，PPO best 未能超过 SFT baseline，说明当前训练分布与测试分布之间仍存在差异。

其次，本文实现的 PPO 框架仍是简化版本。当前方法将一次完整代码生成视为单步 episode，虽然包含 policy loss、value loss、KL penalty 和单步 advantage，但尚未实现完整多步 trajectory PPO。对于语言模型生成任务而言，代码是逐 token 生成的长序列，完整建模每一步 token 的动作、状态和回报传播会更加复杂。本文的单步近似降低了实现难度，也适合有限算力条件下的实验，但其优化能力和稳定性仍弱于完整大规模 PPO 训练框架。

再次，reward 设计仍存在进一步优化空间。reward\_v2 在验证集上取得了最好结果，但它仍主要依赖当前测试用例的通过情况。如果测试用例覆盖不足，模型可能学到通过当前测试的局部策略，而不是获得真正稳健的程序语义理解。reward\_v3 和 reward\_v3b 尝试加入 AST 相似度和函数签名匹配，但实验表明简单结构奖励并未提升 Pass@1。这说明如何设计既密集又与功能正确性高度一致的 reward，仍然是代码生成强化学习中的关键问题。

此外，当前执行环境相对简化。本文采用本地 Python 执行环境完成生成代码测试，能够满足 MBPP 和 HumanEval 的基本实验需求，但还不具备完整沙箱系统的资源隔离和超时控制能力。对于更复杂的代码生成任务，如果生成代码包含死循环、文件操作、网络请求或高资源消耗行为，简单本地执行方式可能存在安全性和稳定性问题。后续若扩展到更复杂数据集，需要引入更严格的沙箱执行机制。

最后，本文主要采用 Pass@1 和 average reward 作为核心评价指标。虽然这些指标能够反映代码是否通过测试，但对代码可读性、复杂度、鲁棒性和安全性等方面的评价仍不充分。实际软件工程场景中的代码质量不仅取决于是否通过当前测试，还涉及可维护性、异常处理、性能和安全性等因素。因此，本文结果更适合说明执行反馈对函数级代码生成正确性的影响，而不能完全代表实际工程代码质量。

\section{研究展望}

针对上述局限，后续研究可以从以下几个方向继续推进。

第一，扩大训练数据规模并提高数据多样性。后续可以引入更多代码生成数据集，如扩展版 MBPP、EvalPlus、APPS 或真实开源项目中的函数级任务，使模型在更丰富的问题类型和测试分布上进行训练。更大规模和更多样化的 rollout 数据有助于缓解当前 PPO 只在验证集分布上提升的问题，并可能提升模型在 HumanEval 等测试集上的泛化能力。

第二，改进 reward 设计。本文实验表明，reward\_v2 比原始 reward 和结构 reward 更有效，但仍然依赖测试用例覆盖。后续可以考虑结合多种反馈信号，例如隐藏测试生成、错误类型分类、静态分析、代码复杂度约束、边界条件覆盖率等，构建更稳健的 reward。同时，也可以研究分阶段 reward，将语法正确性、入口函数一致性、部分测试通过和边界测试通过分别建模，降低奖励噪声。

第三，完善 PPO 训练框架。后续可以实现更完整的多步 PPO，将 token 级生成过程纳入轨迹建模，并引入更稳定的 value function 训练、GAE 估计和 KL 自适应控制。对于更大模型和更大数据集，还可以使用分布式 rollout 和批量执行环境，提高训练效率和样本利用率。

第四，增强执行环境的安全性和鲁棒性。代码生成模型可能产生不可控代码，因此执行反馈训练需要安全隔离。后续可以将本地执行器替换为具备超时控制、资源限制和文件系统隔离能力的沙箱环境，避免死循环或危险操作影响训练过程。同时，可以记录更细粒度的错误类型，为 reward 设计和案例分析提供更多信息。

第五，扩展评价指标和案例分析。除 Pass@1 外，后续可以引入 Pass@k、平均测试通过率、代码复杂度、运行时间和更严格的隐藏测试集评价。对于模型退化案例，可以进一步分析 PPO 更新前后的输出分布变化，研究哪些类型的问题更容易被执行反馈改善，哪些类型的问题容易发生边界条件退化。

总体而言，基于执行反馈的强化学习为提升代码生成模型的功能正确性提供了一条可行路径。本文在小规模数据和有限算力条件下验证了该路径的基本可行性，也揭示了 reward 设计、训练稳定性和泛化能力方面的挑战。未来随着更高质量数据、更完善执行环境和更稳定强化学习算法的发展，执行反馈有望在代码生成模型后训练中发挥更重要作用。

}

%TC:ignore

\clearpage
{
\ExplSyntaxOn
\bool_if:NTF \g__sjtu_twoside_bool
{
    \fancyhead [ LE ]     { 参考文献 }
    \fancyhead [ RO ]     { 参考文献 }
}
{
    \fancyhead [ R ] { 参考文献 }
}
\ExplSyntaxOff
% 文献表字体
\renewcommand{\bibfont}{\zihao{5}}
% 设定固定间距
\fixedlineskip{15.6bp}
{
\ctexset{chapter={afterskip=26bp}}
% 参考文献
\printbibliography[heading=bibintoc]
}
\clearpage
}

\makeatletter
% \appendix采用数字编号。
\renewcommand{\appendix}{\par
    \setcounter{chapter}{0}
    \setcounter{section}{0}
    \ctexset{chapter/number={\arabic{chapter}}}
}
% 使用 \appchapter 替代附录中的 \chapter 章节，附录中的章节不再放入目录。
\newcommand{\appchapter}[1]{
    \refstepcounter{chapter}
    \SJTU@head*[附录 \thechapter]{#1（附录 \thechapter）}
}
\makeatother

{
\ctexset{chapter={afterskip=26bp}}
% 附录
\appendix
% 附录中图表不加入索引
\captionsetup{list=no}
\input{contents/appendix_1}
}


% 结尾部分
\backmatter

{
\ctexset{chapter={afterskip=26bp}}
% 发表论文及科研成果
\input{contents/achievements}
}

\clearpage
{
\ExplSyntaxOn
\bool_if:NTF \g__sjtu_twoside_bool
{
    \fancyhead [ LE ]     { 致谢 }
    \fancyhead [ RO ]     { 致谢 }
}
{
    \fancyhead [ R ] { 致谢 }
}
\ExplSyntaxOff
{
\ctexset{chapter={afterskip=26bp}}
% 致谢
\input{contents/acknowledgements}
}

\clearpage
}

{
\ctexset{chapter={afterskip=26bp}}
% 学士学位论文要求在最后有一个大摘要，单独编页码
% !TEX root = ../main.tex

\begin{digest}
\noindent\textbf{Keywords:} code generation; execution feedback; reinforcement learning; supervised fine-tuning; PPO

\vspace{1em}

Code generation based on natural language instructions has become an important research topic in both artificial intelligence and software engineering. With the rapid development of large language models, modern code-oriented foundation models are able to read problem descriptions, infer algorithmic intentions, and generate executable programs in languages such as Python. These models are usually pretrained on massive code and text corpora, and then adapted to downstream tasks through supervised fine-tuning. This training paradigm has significantly improved the usability of code generation systems. However, code generation is different from ordinary text generation. A fluent and syntactically plausible answer is not necessarily a correct program. The final quality of generated code is determined by whether the program can be executed and whether it satisfies the functional requirements expressed by test cases. Therefore, there is a natural gap between token-level supervised learning objectives and execution-level correctness objectives.

This thesis studies reinforcement learning for code generation based on execution feedback. The central idea is to use the executability of code as a source of training signal. Instead of only asking the model to imitate a reference solution, the proposed workflow lets the model generate candidate programs, runs these programs against unit tests, converts the execution result into reward, and then uses reinforcement learning to optimize the generation policy. In this way, the training objective is more closely aligned with the evaluation criterion of code generation tasks. The work focuses on a small-scale but complete experimental setting: data processing, supervised fine-tuning, rollout collection, local execution, reward computation, advantage estimation, PPO policy update, and final evaluation are all implemented and connected into a full experimental loop.

The motivation of this work comes from two observations. First, a programming problem often has multiple correct solutions. Supervised fine-tuning on a single reference solution may encourage the model to reproduce a particular surface form, but it does not directly teach the model to satisfy the semantic constraints of the task. Second, execution feedback is relatively objective and easy to obtain for many code generation benchmarks. If a generated function passes all provided tests, it is more likely to be functionally correct; if it fails some tests, the number of passed tests and the error type can still provide useful partial feedback. This property makes code generation a suitable domain for reinforcement learning with automatic rewards.

The base model used in this thesis is Qwen2.5-Coder-1.5B, a code-oriented causal language model with strong pretrained programming ability. The experimental data are constructed from MBPP and HumanEval. MBPP is used for supervised fine-tuning, validation, and PPO rollout experiments, while HumanEval is used as the final cross-dataset test set. The raw datasets are converted into two unified formats. The SFT format contains problem prompts and reference responses, and is used for supervised fine-tuning. The RL format keeps the problem prompt, canonical solution, entry point, and test code, and is used for generation, execution, reward calculation, and policy optimization. After processing, the final dataset contains 385 MBPP training samples, 42 MBPP validation samples, and 164 HumanEval test samples.

Before training and evaluation, this thesis performs runtime validation on the processed datasets. For each sample, the canonical solution is executed together with the corresponding tests to ensure that the entry point, solution code, and test code are compatible with the local execution environment. This step is necessary because the reward signal in the PPO stage directly depends on test execution. If the test itself is malformed, or if the entry point cannot be resolved, the calculated reward would be unreliable. The runtime validation results show that all 591 processed samples can be correctly executed in the current framework. This provides a reliable data foundation for the following experiments.

The training pipeline starts with supervised fine-tuning. LoRA is used as the parameter-efficient fine-tuning method, so that the base model can be adapted under limited GPU memory. The SFT stage trains the model on 385 MBPP training samples for five epochs. The target modules include the major attention and feed-forward projection layers. The resulting SFT adapter is used as the baseline policy for subsequent PPO experiments. The SFT validation loss is 0.7472. Under the unified evaluation pipeline, the SFT baseline achieves a Pass@1 of 0.5714 on the MBPP validation set and 0.5244 on the HumanEval test set.

The core contribution of this thesis is the construction of an execution-feedback PPO loop for code generation. In the rollout stage, the current policy model receives a problem prompt and generates a candidate Python solution. Since language models may output explanations, Markdown code fences, repeated problem statements, or additional examples, the generated text must be cleaned before execution. The implemented post-processing module extracts the most plausible code block, aligns the function name with the required entry point when possible, and removes irrelevant text. The cleaned code is then executed in a local Python environment together with the test cases. The execution result includes whether the code is executable, how many tests are passed, whether all tests are passed, and whether an exception or assertion failure occurs.

Based on the execution result, the framework computes a scalar reward for each generated solution. The reward is then normalized into an advantage signal and used by PPO to update the LoRA adapter. The PPO implementation in this thesis is a simplified single-step version. Each generated program is treated as one episode, and the whole response is regarded as one action. The update objective includes the clipped policy loss, a scalar value loss, and a KL penalty against the reference policy. Although this is not a full token-level multi-step PPO implementation, it is sufficient for validating whether execution feedback can improve the generation behavior under limited data and computing resources.

Reward design is one of the key issues in this work. The thesis compares several reward functions. The original reward mainly uses the ratio of passed tests as the reward signal. This design is simple and directly related to functional correctness, but it may still be sparse when most generated solutions fail all tests. To provide denser feedback, \texttt{reward\_v2} combines three components: an executability reward, a partial test pass reward, and a full-pass bonus. This design distinguishes syntax or runtime failure, partial correctness, and complete correctness. It gives the model useful feedback even when the generated solution does not pass all tests.

This thesis also explores more complex reward designs, namely \texttt{reward\_v3} and \texttt{reward\_v3b}. These two variants introduce structural signals such as AST similarity and function signature matching. The intention is to check whether structural similarity to the reference solution can provide additional guidance beyond test execution. However, the experiments show that these structural rewards do not improve Pass@1 in the current setting. In some cases, they increase the average reward value but fail to improve the number of fully correct solutions. This indicates that, for small-scale code generation PPO, a more complex proxy reward may not be better if it is not well aligned with functional correctness.

The main validation results compare Base, SFT baseline, and the best PPO model on the MBPP validation set. The Base model achieves a Pass@1 of 0.6667, with 28 out of 42 samples fully passing the tests. The SFT baseline achieves a Pass@1 of 0.5714, with 24 fully passed samples. This result shows that small-scale supervised fine-tuning does not necessarily improve a strong code base model on every distribution. It may adapt the output style to the training data, but it can also damage part of the pretrained model's original capability. After PPO optimization with \texttt{reward\_v2}, the best PPO model achieves a Pass@1 of 0.6667, also with 28 fully passed samples. Therefore, PPO recovers the validation performance drop caused by SFT and matches the Base model on Pass@1.

The average reward results provide additional information. On the MBPP validation set, Base obtains an average reward of 0.7381, SFT obtains 0.6524, and PPO best obtains 0.7190. Although PPO best has the same Pass@1 as Base, its average reward is slightly lower. This suggests that Base may receive higher partial rewards on some failed samples, while PPO best more directly improves the number of fully passed samples. For code generation, Pass@1 is usually the more important metric because it measures whether a generated answer fully solves the problem in one attempt. Therefore, the validation result supports the conclusion that execution-feedback PPO can correct part of the SFT-induced performance degradation.

To test cross-dataset generalization, this thesis evaluates Base, SFT baseline, and PPO best on HumanEval. The Base model achieves a Pass@1 of 0.5000, the SFT baseline achieves 0.5244, and PPO best achieves 0.5122. The SFT model slightly improves over Base on HumanEval, while PPO best is between Base and SFT. The number of generation execution failures decreases from 7 for Base and 2 for SFT to 0 for PPO best, which shows that PPO improves the execution stability of generated outputs. However, PPO does not outperform SFT on HumanEval Pass@1. This result indicates that the PPO improvement mainly appears on the validation distribution used for rollout and tuning, and it does not fully transfer to the cross-dataset test setting.

The PPO hyperparameter experiments further reveal the sensitivity of reinforcement learning for code generation. The best original-reward PPO setting uses 42 rollouts, a learning rate of $1\times10^{-5}$, a KL coefficient of 0.02, and a batch size of 4. Reducing the KL coefficient to 0.01 leads to worse performance, because the policy moves too far away from the SFT policy and the mean probability ratio becomes too large. Reducing the learning rate to $5\times10^{-6}$ makes the update more conservative, but also weakens the improvement. Increasing the rollout number from 42 to 100 does not automatically improve performance either. These results suggest that PPO effectiveness depends on the balance among reward density, rollout scale, KL constraint, and update strength. More rollout data or a more aggressive update does not necessarily lead to better results.

The reward ablation experiments show that \texttt{reward\_v2} is the most effective design in this thesis. With \texttt{reward\_v2}, PPO best reaches a Pass@1 of 0.6667 on the MBPP validation set. In contrast, PPO with the original reward reaches 0.6190, while \texttt{reward\_v3} and \texttt{reward\_v3b} both reach 0.5952. The structural rewards in \texttt{reward\_v3} and \texttt{reward\_v3b} fail to improve Pass@1 even though they provide additional reward components. This finding is important because it shows that reward functions should not only be dense, but also aligned with the final evaluation target. For code generation, test execution remains the most reliable automatic signal in this experimental setting.

The case analysis provides a more concrete explanation of the model behavior. In a positive case from HumanEval, the task requires sorting the values at even indices of a list while keeping odd indices unchanged. The SFT output incorrectly uses an undefined variable, causing execution failure. The PPO output correctly extracts the even-indexed elements, sorts them, and writes them back to the corresponding positions. This case shows that execution feedback can help correct indexing and implementation logic errors. In another positive case, the task requires extracting fruit counts from a string. The SFT output uses a brittle hard-coded character access pattern, while the PPO output parses numeric tokens from the string and subtracts their sum from the total fruit count. This improves robustness for multi-digit quantities.

There are also negative cases. For example, in a HumanEval task that asks for the product of odd digits and requires returning 0 when all digits are even, Base and SFT correctly maintain a flag to handle the all-even case. PPO generates a shorter implementation based on reduction over odd digits. This implementation works when odd digits exist, but returns 1 instead of 0 when the input contains only even digits. This case shows that PPO can damage boundary condition handling. The reason is that the current PPO update is based on limited rollout data and simplified single-step credit assignment. If the reward does not sufficiently cover certain edge cases, the policy may learn a shortcut that passes many tests but fails important corner cases.

Overall, the experimental results support several conclusions. First, the complete SFT and PPO pipeline for execution-feedback code generation is feasible. The framework can process datasets, fine-tune a base model, generate code, execute tests, compute rewards, and update a policy adapter. Second, execution feedback can help repair the performance drop caused by small-scale SFT on the MBPP validation distribution. Third, reward design is critical. A direct and execution-aligned dense reward such as \texttt{reward\_v2} is more effective than more complicated structural rewards in the current setting. Fourth, PPO is sensitive to KL constraint, learning rate, rollout scale, and reward density. Fifth, the generalization improvement is still limited. PPO best improves validation Pass@1 relative to SFT, but does not surpass SFT on HumanEval.

This thesis also has several limitations. The most important limitation is data scale. The SFT training set contains only 385 MBPP samples, and the main PPO rollout set contains 42 validation samples in the best setting. Such a small scale is enough to verify the experimental pipeline, but it is insufficient for robust policy optimization. The second limitation is the simplified PPO formulation. Treating an entire generated program as one action greatly reduces implementation complexity, but it does not fully model token-level decisions and long-range credit assignment. The third limitation is reward coverage. The reward is based on available tests, so the model may overfit to the current test distribution if the tests are not comprehensive. The fourth limitation is the execution environment. The current local executor is suitable for MBPP and HumanEval style function-level tasks, but a stronger sandbox with timeout, resource isolation, and file system restrictions is required for more complex code generation scenarios.

Future work can improve this research in several directions. First, larger and more diverse training data should be introduced, including extended MBPP variants, EvalPlus, APPS, and real-world repository tasks. More diverse tasks may improve the cross-dataset generalization of PPO. Second, reward design can be improved by combining execution results with hidden tests, error classification, static analysis, boundary condition coverage, and code complexity constraints. Third, the PPO framework can be extended to token-level or multi-step trajectory optimization with a stronger value function and adaptive KL control. Fourth, the execution environment should be upgraded to a safer sandbox with timeout and resource control, so that generated code with infinite loops or unsafe operations can be handled robustly. Finally, evaluation should go beyond Pass@1 and average reward, and include Pass@k, code robustness, running time, readability, and failure type analysis.

In conclusion, this thesis verifies that execution feedback can be used as an effective reinforcement learning signal for code generation. Under small-scale data and limited computing resources, PPO based on execution feedback can improve the SFT model on the validation distribution and reduce execution failures. At the same time, the experiments reveal that strong code base models already have considerable capability, and that small-scale SFT and PPO do not automatically guarantee better generalization. The most practical lesson of this work is that execution-based reward is valuable, but it must be carefully designed and constrained. A simple, direct, and functionally aligned reward can be more effective than a complex proxy reward. This provides a concrete experimental reference for future research on reinforcement learning and execution-based optimization in code generation.
\end{digest}

}

\end{document}
